%%%%%%%%%%%%%%%%%%%%%%%%%%%%%%%%%%%%%%%%%%%%%%%%%%%%%%%%%%%%%%%%%%%%%%%%%
%                                                                       %
%                        Chapter 5 - Section 1                          %
%                                                                       %
%%%%%%%%%%%%%%%%%%%%%%%%%%%%%%%%%%%%%%%%%%%%%%%%%%%%%%%%%%%%%%%%%%%%%%%%%     


\chapter{Techniques of Differentiation}

In this chapter we focus on functions given by {\em formulas\/}.  The
derivatives of such functions are then also given by formulas.  In
chapter 4 we used information about the derivative of a function to
recover the function itself; now we go {\em from} the function {\em
  to} its derivative.  We develop the rules for {\em differentiating}
a function: computing the formula for its derivative from the formula
for the function.  Then we use differentiation to investigate the
properties of functions, especially their {\em extreme
  values}. Finally we examine a powerful method for solving equations
that depends on being able to find a formula for a derivative.

\section{The Differentiation Rules}

There are three kinds of differentiation rules.  First, any basic
function has a specific rule giving its derivative.  Second, the {\em
  chain rule} will find the derivative of a {\em chain} of functions.
Third, there are general rules that allow us to calculate the
derivatives of algebraic combinations---e.g., sums, products, and
quotients---of any functions provided we know the derivatives of each
of the component functions. To obtain all three kinds of rules we will
typically start with the analytic definition of the derivative as the
limit of a quotient of differences:

\begin{center}
\begin{picture}(360,20)
\put(0,-46){\framebox (360,68){
\begin{minipage}{330pt}
{\bf Definition}.  The {\bf derivative}\index{derivative} of the function $f$ at $x$ is the value of the limit
$$
\lim_{\Delta x \rightarrow 0} \frac{f(x + \Delta x) - f(x)}{\Delta x} = f'(x).
$$
\end{minipage}
}}
\end{picture}
\end{center}

\newpage

In this chapter we will look at the cases where this limit can be
evaluated exactly.  Although using this definition of derivative
usually leads to many algebraic manipulations, the other
interpretations of derivatives as slopes, rates, and multipliers will
still be helpful in visualizing what's going on.  The process of
calculating the derivative of a function is called \mbox{\bf
  differentiation}\index{differentiation}.  For this reason, functions
which are locally linear and not locally vertical (so they do have
slopes, and hence derivatives at every point) are called {\bf
  differentiable}\index{differentiable} functions.  Our goal in this
chapter is to differentiate functions given by formulas.

 

\subsection{Derivatives of Basic Functions}

      When a function is given by a {\em formula\/},  
      \mar{Functions given by formulas have derivatives given by formulas}
there is in fact a formula for its derivative. We have already seen several examples in chapters 3 and 4.  These examples include all of what we may consider the {\bf basic functions}.  We collect these formulas in the following table.
\begin{center}
{\bf Rules for Derivatives of Basic Functions}
$$
\begin{array}{cccc}
   \mbox{function}    &\qquad &        &\mbox{derivative} \\[2pt]  
   \hline \\ [-9pt]
        mx+b            & &      &m \\
        x^{r}           & &              &rx^{r-1} \\
        \sin{x}         & &              &\cos{x} \\
        \cos{x}         & &              &-\sin{x}  \\
        e^{x}           & &              &e^{x} \\
        \ln{x}          & &              &1/x \\

\end{array}
$$
\end{center}

In the case of the linear function $mx+b$, we obtained the derivative
by using its geometric description as the {\em slope\/} of the graph
of the function.  The derivatives of the exponential and logarithm
functions came from the definition of the exponential function as the
solution of an initial value problem.  To find the derivatives of the
other functions we will need to start from the definition.
	
\subsubsection{An example: $f(x)=x^3$} 

We begin by examining the calculation of the derivative of $f(x) =
x^3$ using the definition. The change $\Delta y$ in $y = f(x)$
corresponding to a change $\Delta x$ in $x$ is given by
	
\mbox{}
\vspace{-35pt}

\begin{align*}
   \Delta y &= f(x + \Delta x) - f(x) \\
            &= (x + \Delta x)^{3} - x^{3} \\
            &= 3x^{2}\cdot\Delta x + 3x(\Delta x)^{2} + (\Delta x)^{3}. 
\end{align*}
From this we get 
\begin{align*}
f'(x) &= \lim_{\Delta x \rightarrow 0} \frac{\Delta y}{\Delta x} \\[3pt]
      &= \lim_{\Delta x \rightarrow 0} 3x^{2} + 3x\cdot\Delta x + 
      	(\Delta x)^{2}.
\end{align*}

To see what's happening with this expression, let's consider the
specific value $x=2$ and evaluate the corresponding values of $\Delta
y/ \Delta x\/$ for successively smaller
%
\mar{\vspace{40pt}The value of $\Delta y/ \Delta x$ gets closer and
  closer to 12 as $\Delta x$ gets smaller and smaller}
%
$\Delta x\/$.

\begin{center}
\begin{tabular}{lll}
$\Delta x$ & $2^2 + 6\Delta x + (\Delta x)^2$ & $\Delta y/\Delta x$ \\[2pt] \hline \\ [-8pt]
.1 & 12 + .6 + .01 & 12.61 \\
.01 & 12 + .06 + .0001 & 12.0601 \\
.001 & 12 + .006 + .000001 & 12.006001 \\
.0001 & 12 + .0006 + .00000001 & 12.00060001 \\
.00001 & 12 + .00006 + .0000000001 & 12.0000600001
\end{tabular}
\end{center}
It is clear from this table that we can make $\Delta y$/$\Delta x$ as
close to 12 as we like by making $\Delta x$ small enough.  Therefore
$f'(2)=12$.

Note that in the table above we have used positive values of $\Delta
x$.  You should check to convince yourself that if we had used
negative values of $\Delta x$ we would have come up with a different
set of approximations $\Delta y$/$\Delta x$, but that the limit would
still be the same, namely 12---it doesn't matter whether we use
positive or negative values for $\Delta x$, or a mixture of the two,
so long as $\Delta x \rightarrow 0$.

In general, for any given $x$, the second and third terms in the
expansion for $\Delta y$/$\Delta x$ become vanishingly small as
$\Delta x \rightarrow 0$, so that $\Delta y$/$\Delta x$ can be made as
close to $3x^{2}$ as we like by making $\Delta x$ small enough.  For
this reason, we say that the derivative $f'(x)$ is {\em exactly\/}
$3x^{2}$ :
	$$
	f'(x) = \lim_{\Delta x \rightarrow 0} 3x^2 + 3x\cdot\Delta x + 
		(\Delta x)^2 = 3x^2.
	$$
In other words, given the function $f$ specified by the formula $f(x)=x^3$ we have found the formula for its derivative function $f'$: $f'(x) = 3 x^2$.  Note that this general formula agrees with the specific value $f'(2)=12$ we have already obtained.
\bigskip

	Notice the difference between the statements 
	$$
	f'(x) \approx \Delta y/\Delta x \qquad \mbox{and} \qquad
		f'(x) = 3x^2.
	$$
For a particular value of $\Delta x\/$, the corresponding value of $\Delta y/\Delta x$ is an approximation of $f'(x)\/$.  We can obtain another, better approximation by computing $\Delta y/\Delta x$ for a smaller $\Delta x\/$.  The successively better approximations differ from one another by less and less.  In particular, they differ less and less from the {\em limit value} $3x^2$.  The value of the derivative $f'(x)$ is {\em exactly} $3x^2\/$.

	More generally, for any function $y=f(x)$, a particular difference quotient $\Delta y/\Delta x$ is an approximation of $f'(x)\/$.  Successively smaller values of $\Delta x$ give successively better approximations of $f'(x)\/$.  Again $f'(x)$ {\em exactly} equals the limiting value of these successive approximations.   In some cases, however, we are only able to approximate that limiting value, as we often did in chapter 3, and for many purposes the approximation is entirely satisfactory.  In this chapter we will concentrate on the exact statements that are possible for functions given by formulas.

\subsubsection{The other basic functions} 

	Our formula for the derivative of the function  $f(x) = x^{3}$ is one instance of the general rule for the derivative 
	\mar{\vspace{22pt}The rule for\\ the derivative of\\ a power function}
of $f(x) = x^r\/$.
\begin{center}
\begin{picture}(300,50)
	\put(0,0) 	{\framebox(300,50)
	{\shortstack
	{	
	\bf For every real number \boldmath $r$ \unboldmath, 			the derivative   \\
	\bf of \hspace{.1in} \boldmath $f(x) = x^{\displaystyle r}$ \unboldmath \hspace{.1in} 	      is \hspace{.1in} \boldmath $f'(x) = r\, x^{\displaystyle r-1}$. \unboldmath
	}}}
\end{picture}
\end{center}

We can prove this rule for the case when $r$ is a positive integer
using algebraic manipulations very like the ones carried out for
$x^{3}$; see the exercises for verifications of this and the other
differentiation rules in this section. Using a rule for quotients of
functions (coming later in this section), we can show that this rule
also holds for {\em negative\/} integer exponents. Further arguments
using the chain rule show that the pattern still holds for {\em
  rational\/} exponents. We can eliminate this case-by-case approach,
though, by recalling the approach developed in chapter 4.  We saw that
we can give meaning to $b^{r}$ for any positive base $b$ and any real
number $r$ by defining
$$
b^{\textstyle r} = e^{\textstyle r \ln(b)}.
$$
Using the formulas for the derivatives of $e^{\displaystyle x}$ and ln $x$ together with the chain rule, we can prove the rule for $x>0$ and for arbitrary real exponent $r$ directly, without first proving the special cases for integer or rational exponents.  See the exercises for details.  Arguments justifying the formulas for the derivatives of the trigonometric functions are also in the exercises.

\subsection{Combining Functions}

	We can form new functions by combining functions.  We have already studied one of the most useful ways of doing this in chapter 3 when we looked at forming ``chains'' of functions and 
	\mar{Functions combined\\ by chains\ldots }	
developed the {\bf chain rule} for taking the derivative of such a chain.  Suppose $u=f(x)$ and $y=g(u)$.  Chaining these two functions together we have $y$ as a function of $x$:
$$ 
y = h(x) = g(f(x)).
$$  
The chain rule tells us how to find the derivative of $y$ with respect to $x$.  In function notation it takes the form
$$
h'(x) = g'(f(x)) \cdot f'(x).
$$ 
In Leibniz notation, using $f(x) = u$ we can write the chain rule  
	\mar{\vspace{21pt}The chain rule}
as
$$
\frac{dy}{dx} = \frac{dy}{du} \cdot \frac{du}{dx}.
$$

	We also saw in chapter 3 that the polynomial $5x^{3}-7x^2 +3$ can be thought of as an {\em algebraic\/} combination of simple functions.   
	\mar{\ldots and algebraically}	
We can build an even more complicated function by forming a quotient with this polynomial in the numerator and the difference of the functions $\sin{x}$ and $e^x$ in the denominator.  The result is
$$
\frac{5x^{3}-7x^2 +3}{\sin{x} - e^x}.
$$

	   The derivative of this function, as well as of other functions formed by adding, subtracting, multiplying and dividing simpler functions, is  obtained by use of the following rules for the derivatives of algebraic combinations of differentiable functions.

\newpage

\begin{center}
{\bf Rules for Algebraic Combinations of Functions}
\mar{\vspace{50pt}Combining functions by adding, subtracting, multiplying and dividing}
$$
\begin{array}{cccc}
   \mbox{function}    & &        &\mbox{derivative} \\[.05in]  \hline

  \mbox{\rule{0in}{.25in}$f(x) + g(x)$}           & & &\mbox{\rule{0in}{.25in}$f'(x) + g'(x)$}       \\[.1in]
        f(x) - g(x)     & &      &f'(x) - g'(x)  \\[.1in]
        cf(x)           & &      &cf'(x)  \\[.1in]
        f(x)\cdot g(x)     &  &       &f'(x)\cdot g(x) + f(x)\cdot g'(x)  \\[.1in]
\displaystyle{\frac{f(x)}{g(x)}}   & &
&\displaystyle{ \frac{g(x)\cdot f'(x)-f(x)\cdot g'(x)}{[g(x)]^2} }
\end{array}
$$
\end{center}
Notice carefully that the product rule has a plus sign but the quotient rule has a {\em minus\/} sign.  
	\mar{Notice the signs\\ in the rules}
You can remember these formulas better if you think about where these signs come from.  Increasing {\em either\/} factor increases a (positive) product, so the derivative of {\em each\/} factor  appears with a plus sign in the formula for the derivative of a product.  Similarly, increasing the numerator increases a positive quotient, so the derivative of the numerator appears with a plus sign in the formula for the derivative of a quotient.  However, increasing the denominator {\em decreases\/} a positive quotient, so the derivative of the denominator appears with a minus sign.

	Let's now use the rules to differentiate the quotient
	$$
	\dfrac{5x^3-7x^2+3}{\sin{x} -e^x}.
	$$
First, the derivative of the numerator $5x^{3}-7x^{2}+3$ is 
$$
5(3x^{2}) - 7(2x) + 0 = 15x^{2} - 14x.  
$$
Similarly the derivative of $\sin{x} - e^x$ is $\cos{x} - e^x$.  Finally, the derivative of the quotient function is obtained by using the rule for quotients:
	$$
\frac{(\sin{x} - e^x)(15x^{2} - 14x) - (5x^{3} - 7x^{2} + 3)(\cos{x} - e^x)}
{(\sin{x} - e^x)^{2}}.
	$$

\newpage

	The following examples further illustrate the use of the rules for algebraic combinations of functions.
$$
\begin{array}{cccc}
   \mbox{function}    &  &        &\mbox{derivative} \\[.05in]  \hline
\mbox{\rule{0in}{.25in}$-3e^{t} + \sqrt[3]{t}$}    &  &        &\mbox{\rule{0in}{.25in}$-3e^{t} + (1/3)t^{-2/3}$}  \\[.1in]
\displaystyle{\frac{5}{x^3}} -7x^4+\ln x & & &5(-3)x^{-4} - 7(4x^3) + 1/x  \\[.1in]
\displaystyle{7 \sqrt{x} \cos x}     & &      &7 \displaystyle{(\frac{1}{2\sqrt{x}})} \cos x +7\sqrt{x}(-\sin x) \\[.1in]
\displaystyle{\left(\frac43\right)} \pi r^3  &     &   &\displaystyle{\left(\frac43\right) \pi 3r^2} \\[.22in]
\displaystyle{ \frac{3s^6}{s^2 - s} }   &   &     &\displaystyle{ \frac{(s^2-s)3(6s^5) -3s^6(2s-1)}{(s^{2} - s)^2} }   \end{array}
$$

	For another kind of example, suppose the per capita daily energy consumption in a country is currently 800,000 BTU, and, due to energy conservation efforts, it is falling at the rate of 1,000 BTU per year.  
Suppose too that the population of the country is currently 200,000,000 people and is rising at the rate of 1,000,000 people per year.  Is the total daily energy consumption of this country rising or falling?  By how much?

Three different quantities vary with time in this example:  daily per capita energy consumption, population and total daily energy consumption.  We can model this situation with three functions  $C(t)$, $P(t)$ and $E(t)$.
\begin{align*}
C(t) &: \text{per capita consumption at time $t$} \\
P(t) &: \text{population at time $t$} \\
E(t) &: \text{total energy consumption at time $t$}
\end{align*}

Since the per capita consumption times the number of people in the
population gives the total energy consumption, these three functions
are related algebraically:
$$
E(t) = C(t) \cdot P(t).
$$
If $t=0$ represents today, then we are given the two rates of change
\begin{align*}
C'(0) & =  -1,000 = - 10^{3} \, \mbox{BTU per person per year, and} \\
P'(0) & =  1,000,000 = 10^{6} \, \mbox{persons per year.}
\end{align*}


Using the product rule we can compute the current rate of change of
the total daily energy consumption:

\newpage

\mbox{}

\vspace{-30pt}

\begin{align*}
E'(0) &= C(0) \cdot P'(0) + C'(0) \cdot P(0) \\
      &= (8 \times 10^{5})\cdot (10^{6}) + (-10^{3})\cdot(2 \times 10^{8}) \\
      &= (8 \times 10^{11}) - (2 \times 10^{11}) \\
      &= 6 \times 10^{11} \, \mbox{BTU per year}.
\end{align*}
So the total daily energy consumption is currently rising at the rate
of $6 \times 10^{11}$ BTU per year.  Thus the growth in the population
more than offsets the efforts to conserve energy.

	Finally, it is a useful exercise to check that the units make sense in this computation.  
	\mar{Checking units}	
Recall that $C(t)$ represents {\em per capita} daily energy consumption, so the units for $C(0) \cdot P'(0)$ are
$$
\frac{\mbox {BTU}}{\mbox {person}} \cdot \frac{\mbox{persons}}{\mbox{year}} = \frac{\mbox{BTU}}{\mbox{year}},
$$
and, similarly, the units for $C'(0) \cdot P(0)$ are
$$
\frac{\mbox{BTU}}{\mbox{person}} \cdot \frac{1}{\mbox{year}} \cdot 
{\mbox {persons}} = \frac{\mbox{BTU}}{\mbox{year}}.
$$


\subsection{Informal Arguments}

	All of the rules for differentiating algebraic combinations of functions can be proved by using the algebraic definition of the derivative as a limit of a difference quotient.  In fact, we will examine such a formal proof below.  However, informal arguments based on geometric ideas or other intuitive insights are also valuable aids to understanding.  Here are three examples of such arguments.

\begin{itemize}

\enlargethispage*{60pt}
\item If a new function $g$ is obtained from $f$ by multiplying  by a positive constant $c$, so $g(x)=cf(x)$, what is the relationship between the graphs of  $y=f(x)$  and of $y=g(x)$?   
	\mar{Stretching $y$-coordinates}
Stretching (or compressing, if $c$ is less than 1) the $y$-coordinates of the points of the graph of  $f$ by a factor of $c$ yields the graph of $g$.

\vspace{120pt}

\special{psfile=../ps/calc1/pic200a.ps}

\noindent 
What then is the relationship between the slopes $f'(x)$ and $g'(x)$?
If the $y$-coordinates are tripled, the slope will be three times as
great.  If they are halved, the slope will also be half as much.  More
generally, the elongated (or compressed) graph of $g$ has a slope
equal to $c$ times the slope of the original graph of $f$.  In other
words, $g(x)=cf(x)$ implies $g'(x)=cf'(x)$.

\item Now suppose instead that $g$ is obtained from $f$ by adding a
  constant $b$, so $g(x) = f(x) + b$.  
%
\mar{Shifting $y$-coordinates}
%
This time the graph of $y=g(x)$ is obtained from the graph of $y=f(x)$
by shifting up or down (according to the sign of $c$) by $|c|$
units. What is the relationship between the slopes $f'(x)$ and
$g'(x)$?  The shifted graph has exactly the same slope as the original
graph, so in this case, $g(x)=f(x)+b$ implies $g'(x)=f'(x)$.

\vspace{120pt}

\special{psfile=../ps/calc1/pic200b.ps}


\end{itemize}

There is a similar pattern when the coordinates of the input variable
are stretched or shifted---that is when $y=f(u)$ and $u$ is {\em
  rescaled\/} by the linear relation $u=mx+b$.  These results depend
on the chain rule and appear in the exercises.
  
The fact that the derivative of $f(x)+b$ is the same as the derivative
of $f(x)$ is a special case of the general {\em addition rule\/},
which says {\em the derivative of a sum is the sum of the
  derivatives\/}.  In the special case, the derivative of the constant
function $b$ is zero, so adding a constant leaves the derivative
unchanged.  To see that how natural it is to add rates in the general
case, consider the following example:

Suppose we are diluting concentrated orange juice 
%
\mar{Adding flows} 
%
by mixing it with water in a big tub.  We may let $f(t)$ be the amount
(in gallons) of concentrate in the tub and $g(t)$ be the amount of
water in the tub at time $t$.  Then $f'(t)$ is the rate at which
concentrate is being added at time $t$ (measured in gallons per
minute), and $g'(t)$ is the rate at which water is flowing into the
tub.  The formula $F(t) = f(t) + g(t)$ then gives the total amount of
liquid in the tub at time $t$, and $F'(t)$ is the rate by which that
total amount of liquid is changing at time $t$.  Clearly that rate is
the sum of the rates of flow of concentrate and water into the tank.
If at some particular moment we are adding concentrate at the rate of
3.2 gal/min and water at the rate of 1.1 gal/min, the liquid in the
tub is increasing by 4.3 gal/min at that moment.


\subsection{A Formal Proof: the Product Rule}

We include here the algebraic calculations yielding the rule for the
derivative of the product of two arbitrary functions---just to give
the flavor of these arguments.  Algebraic arguments for the rest of
these rules may be found in the exercises.

\medskip

{\bf The Product Rule:}
$$
F(x) = f(x) \cdot g(x) \mbox{ \em implies\/}\ F'(x) = f'(x)\cdot g(x) + f(x)\cdot g'(x)
$$
To save some writing,  let
\begin{align*}
\Delta F &= F(x + \Delta x) - F(x), \\
\Delta f &= f(x + \Delta x) - f(x),   \\
\mbox{and} \qquad \Delta g &= g(x + \Delta x) - g(x).
\end{align*}
Rewrite the last two equations as
\begin{align*}
f(x + \Delta x) &= f(x) + \Delta f  \\
g(x + \Delta x) &= g(x) + \Delta g.
\end{align*}
Now we can write
\begin{align*}
F(x + \Delta x) &= f(x + \Delta x)\cdot g(x + \Delta x)  \\
                &= (f(x) + \Delta f) \cdot (g(x) + \Delta g)  \\
       &= f(x) \cdot g(x) + f(x) \cdot \Delta g + \Delta f \cdot g(x) 
                + \Delta f \cdot \Delta g
\end{align*}
	\mar{Simplifying $\Delta F$}
This gives us a simple expression for
$$
\Delta F = F(x + \Delta x) - F(x)
$$
namely,
$$
\Delta F = f(x) \cdot \Delta g + \Delta f\cdot g(x) + \Delta f\cdot \Delta g
$$

\newpage


These quantities all have nice geometric interpretations.  
%
\mar{Interpret $\Delta f$ and $\Delta g$ as lengths}
%
First, think of the numbers $f(x)$ and $g(x)$ as lengths that depend
on $x$; then $F(x)$ naturally stands for the area of the rectangle
whose sides are $f(x)$ and $g(x)$.  If the sides of the rectangle grow
by the amounts $\Delta f$ and $\Delta g$, then the area $F$ grows by
$\Delta F$.  As the following diagram shows, $\Delta F$ has three
parts, corresponding to the three terms in the expression we derived
algebraically for $\Delta F$.

\begin{center}
\begin{picture}(350,180)(-33,0)
 \put(0,30){\framebox(220,130){area = $f(x)\cdot g(x)$}} % big box
\put(0,0){\framebox(220,30){area = $f(x)\cdot \Delta g$}} % next big
\put(250,15){\vector(-1,0){20}} 		% bottom area stuff
\put(251,13){area = $\Delta f\cdot\Delta g$}
\put(250,95){\vector(-1,0){20}}			% top area stuff
\put(251,93){area = $\Delta f\cdot g(x)$}
\put(220,0){\framebox(20,30){}}		% little box
\put(220,30){\framebox(20,130){}}	% long box
	% bottom g brace
\put(-30,1.5){\makebox(30,27){$\Delta g \left\{ \makebox(0,16){} \right.$}}
	% top g brace
\put(-33,31.5){\makebox(30,130){$g(x) \left\{ \makebox(0,66){} \right.$}}
	% left f brace
\put(2.50,163){$\overbrace{\makebox(215,0){}}^{\textstyle f(x)}$}
	%  right f brace
\put(222.5,163){$\overbrace{\makebox(13,0){}}^{\textstyle \Delta f}$}
\end{picture}
\end{center}


Now we divide  $\Delta F$  by  $\Delta x$  and finish the argument:
\begin{align*}
\frac{\Delta F}{\Delta x} &= \frac{f(x)\cdot \Delta g + \Delta f\cdot g(x) +
\Delta f \cdot \Delta g}{\Delta x} \\ 
     &= f(x) \cdot\frac{\Delta g}{\Delta x}  + \frac{\Delta f}{\Delta x} \cdot g(x) + \frac{\Delta f \cdot \Delta g}{\Delta x}
\end{align*}



Consider what happens to each of the three terms as $\Delta x$ gets
smaller and smaller.  In the first term, the second factor $\Delta g/
\Delta x$ approaches $g'(x)$---by the {\em definition\/} of the
derivative.  The first factor, $f(x)$ doesn't change at all as $\Delta
x$ shrinks.  So the first term approaches $f(x) \cdot g'(x)$.
Similarly, in the second term, the quotient $\Delta f/ \Delta x$
approaches $f'(x)$, and the second term approaches $f'(x) \cdot g(x)$.

Finally, look at the third term.  We would know what to expect if we
had another factor of $\Delta x$ in the denominator.  We can put
ourselves in familiar territory by the ``trick'' of multiplying the
third term by $\Delta x/ \Delta x$:
$$
\frac{\Delta f \cdot \Delta g}{\Delta x} 
= \frac{\Delta f}{\Delta x} \cdot \frac{\Delta g}{\Delta x} \cdot \Delta x
$$ 
Thus we can see that as $\Delta x$ approaches zero, the third term
itself approaches $f'(x) \cdot g'(x) \cdot 0 = 0$.

\bigskip

     We may summarize our calculation by writing
\begin{align*}
\lim_{\Delta x \rightarrow 0} \frac{\Delta F}{\Delta x} &= 
      f(x)\cdot\left(\lim_{\Delta x \rightarrow 0} 
      \frac{\Delta g}{\Delta x}\right)  +  
      \left(\lim_{\Delta x \rightarrow 0} 
      \frac{\Delta f}{\Delta x}\right) \cdot g(x) \\
&    \quad + \left(\lim_{\Delta x \rightarrow 0} 
      \frac{\Delta f}{\Delta x}\right)\cdot 
      \left(\lim_{\Delta x \rightarrow 0}
      \frac{\Delta g}{\Delta x}\right) \cdot
      \left(\lim_{\Delta x \rightarrow 0} \Delta x\right)
\end{align*}
from which we have
\begin{align*}
\lim_{\Delta x \rightarrow 0} \frac{\Delta F}{\Delta x} &= 
   f(x) \cdot g'(x) + f'(x) \cdot g(x) + f'(x) \cdot g'(x) \cdot 0  \\
            &= f(x) \cdot g'(x) + f'(x) \cdot g(x).
\end{align*}
This completes the proof of the product rule.  Other formal arguments are left to the exercises.

\subsection{Exercises}

\subsubsection{Finding Derivatives}

\vspace{-10pt}
\num Find the derivative of each of the following functions.
\medskip

\noindent
\begin {tabular}{ll}
a) $\;\; 3x^5 - 10x^2 + 8$  &  j) $\;\; x^2e^x$ \\
b) $\;\; (5x^{12} + 2)(\pi - {\pi}^2 x^4)$  & k) $\;\; \cos{x} + e^x$ \\
c) $\;\; \sqrt{u} - 3/u^3 + 2u^7$ 
	& l) $\;\;\sin{x}/ \cos{x}$ \\
d) $ \;\; mx + b \;\; (m,b \; \mbox{constant})$ \rule{30pt}{0pt}
	& m) $\;\;e^x \ln{x}$ \\
e) $ \;\; .5 \sin x + \sqrt[3]{x} + \pi^2$ 
	& n) $\;\; \dfrac{2^x}{10 + \sin{x}}$ \\
f) $\;\; \dfrac{{\pi - {\pi}^2 x^4}}{5x^{12} + 2}$  
	& o) $ \;\; \sin(e^x \cos x)$ \\
g) $ \;\;2\sqrt{x} - \dfrac{1}{\sqrt{x}}$ 
	& p) $ \;\; 6 e^{\cos t } / \, 5  \sqrt[3]{t}$ \\
h) $ \;\; \tan z \, (\sin z - 5)$ & q) $ \;\; \mbox{ln}(x^2 + x  e^x)$\\ [2pt]
i) $ \;\; \dfrac{\sin x}{x^2}$ & r) $\;\; \dfrac{5x^2 + \mbox{ln} x}{7 \sqrt{x} + 5}$ 
\end{tabular}

\num Suppose $f$ and $g$ are functions and that we are given
\begin{xalignat*}{3}
f(2)  &= 3, &  g(2) &= 4, & g(3) &= 2,  \\
f'(2) &= 2, &      g'(2) &= -1, &    g'(3) &= 17.
\end{xalignat*}
Evaluate the derivative of each of the following functions at $t=2$:

\smallskip
\noindent
\begin {tabular}{ll}
a) $\;\; f(t) + g(t) $ &  f) $\;\; {\displaystyle \sqrt{g(t)}}$ \\
b) $\;\; 5f(t)-2g(t) $& g) $\;\; t^2f(t)$ \\
c) $\;\; f(t)g(t) $
	& h) $\;\;(f(t))^2+(g(t))^2$ \\ [2pt]
d) $ \;\; \dfrac{f(t)}{g(t)}$ \rule{70pt}{0pt}
	& i) $\;\;\dfrac{1}{f(t)}$ \\ [4pt]
e) $ \;\; g(f(t)$ 
	& j) $\;\;f(3t-(g(1+t))^2)$ \\
\end{tabular}

\smallskip
\noindent
k)\enskip\thinspace What additional piece of information would you
need to calculate the derivative of $f(g(t))$ at $t=2$?

\smallskip
\noindent
l)\enskip\thinspace Estimate the value of $f(t)/g(t)$ at $t=1.95$

\numa Extend the product rule to express $(f(t)g(t)h(t))'$ in terms of
$f$, $g$, and $h$.  \sub If the length, width, and height of a
rectangular box are changing at the rates of 3, 6, and $-5$
inches/minute at the moment when all three dimensions happen to be 10
inches, at what rate is the volume of the box changing then?  \sub If
the length, width, and height of a box are 10 inches, 12 inches, and 8
inches, respectively, and if the length and height of the box are
changing at the rates of 3 inches/minute and $-2$ inches/minute,
respectively, at what rate must the width be changing to keep the
volume of the box constant?


\num In this problem we examine the effect of stretching or shifting
the coordinates of the input variable of a function.  Your answers
should address both the {\em algebra} and the {\em geometry} of the
problem to show how the algebraic relations between the functions are
manifested in their graphs.

\sub Suppose $f(x)=\sin(x)$ and $g(x)= \sin(mx)$, where $m$ is a
constant stretching factor.  What is the relation between $f'(x)$ and
$g'(x)$?

\sub As in (a), suppose $f(x) = \sin(x)$, but this time $g(x) =
\sin(x+b)$ where $b$ is the size of a (constant) shift.  What is the
relation between $f'(x)$ and $g'(x)$ this time?

\sub Now consider the general case: $f(x)$ in an unspecified
differentiable function and $g(x) = f(mx+b)$, where the input variable
is stretched by the constant factor $m$ and shifted by the constant
amount $b$.  What is the relation between $f'(x)$ and $g'(x)$ in this
general case?

\newpage
 
\num Which of the following functions has a derivative which is always
positive (except at $x=0$, where neither the function nor its
derivative is defined)?
	$$
 1/x \qquad    -1/x   \qquad   1/{x^2}  \qquad -1/{x^2}
	$$

\vspace{-8pt}

\numa As a function of its radius $r$, the volume of a sphere is given
by the formula $V(r) = \frac43 \pi r^3$.  If $r$ is measured in
centimeters, what are the units for $V'(r)\/$?

\sub Explain why square~cm are {\em not} the appropriate units for
$V'(r)\/$, even though dimensionally correct.

\num  Do the following. 

\sub Show that $\dfrac{1}{{1 - x^2}}$ and $\dfrac{x^2}{{1 - x^2}}$
have the same derivative.

\sub If $f'(x) = g'(x)$ for every $x$, what can be concluded about the
relationship between $f$ and $g$?  (Hint: What is $(f(x) - g(x))'$ ?)

\sub Show that $\dfrac{1}{1 - x^2} = \dfrac{x^2}{{1 - x^2}} + C$ by
  finding $C$.

\num Suppose that the current total daily energy consumption in a
particular country is $16 \times 10^{13}$ BTU and is rising at the
rate of $6 \times 10^{11}$ BTU per year.  Suppose that the current
population is $2 \times 10^{8}$ people and is rising at the rate of
$10^6$ people per year.  What is the current daily per capita energy
consumption?  Is it rising or falling?  By how much?

\num The population of a particular country is 15,000,000 people and
is growing at the rate of 10,000 people per year.  In the same country
the per capita yearly expenditure for energy is \$1,000 per person and
is growing at the rate of \$8 per year.  What is the country's current
total yearly energy expenditure?  How fast is the country's total
yearly energy expenditure growing?

\num The population of a particular country is 30 million and is
rising at the rate of 4,000 people per year.  The total yearly
personal income in the country is 20 billion dollars, and it is rising
at the rate of 500 million dollars per year.  What is the current per
capita personal income?  Is it rising or falling?  By how much?

\num An explorer is marooned on an iceberg.  The top of the iceberg is
shaped like a square with sides of length 100 feet.  The length of the
sides is shrinking at the rate of two feet per day.  How fast is the
area of the top of the iceberg shrinking?  Assuming the sides continue
to shrink at the rate of two feet per day, what will be the dimensions
of the top of the iceberg in five days?  How fast will the area of the
top of the iceberg be shrinking then?

\num Suppose the iceberg of problem 9 is shaped like a cube.  How fast
is the volume of the cube shrinking when the sides have length 100
feet?  How fast after five days?


\subsubsection{Deriving Differentiation Rules}

\vspace{-10pt}

\num In this problem we calculate the derivative of $f(x) = x^4$.

\sub Expand $f(x + \Delta x) = (x + \Delta x)^4 =$ $(x + \Delta x)(x +
\Delta x)(x + \Delta x)(x + \Delta x)$ as a sum of 16 terms.  (Don't
collect ``like'' terms yet.)

\sub How many terms in part a involve {\em no\/} $\Delta x$'s?  What
form do such terms have?

\sub How many terms in part a involve exactly {\em one\/} $\Delta x$?
What form do such terms have?

\sub  Group the terms in  part a  so that  $f(x + \Delta x)$  has the form
	$$ 
	Ax^4 + B \Delta x + R (\Delta x)^2 \ ,
	$$  
where there are no  $\Delta x$'s  among the terms in  $A$  or  $B$, but  $R$  has several terms, some involving  $\Delta x$.  Use part b to check your value of  $A$;  use part c to check your value of  $B$.
\sub  Compute the quotient  
$\dfrac{f(x + \Delta x) - f(x)}{\Delta x}$, taking advantage of  part d. 
\sub  Now find
	$$
	\lim_{\Delta x \rightarrow 0} 
\dfrac{f(x + \Delta x) - f(x)}{\Delta x}; 
	$$
this is the derivative of $x^4$.  Is your result here compatible with the rule for the derivative of $x^n$ ?

\num In this problem we calculate the derivative of $f(x) = x^n$,
where $n$ is any positive integer.  

\sub First show that you can write
	$$
f(x + \Delta x) = x^n + nx^{n-1} \Delta x + R (\Delta x)^2
	$$
by developing the following line of argument.  Write  $(x + \Delta x)^n$  as a product of  $n$  identical factors:
	$$
(x + \Delta x)^n = \underbrace{(x + \Delta x)}_{1\mbox{-st}}
                   \underbrace{(x + \Delta x)}_{2\mbox{-nd}}
                   \underbrace{(x + \Delta x)}_{3\mbox{-rd}} 
            \ldots \underbrace{(x + \Delta x)}_{n\mbox{-th}}
	$$ 
But now, before tackling this general case, look at the following
examples.  In the examples we use notation to help us keep track of
which factors are contributing to the final result.

\noindent 
i) Consider the product $(a + b)(\underline{a} + \underline{b}) =
a\underline{a} + a\underline{b} + b\underline{a} + b\underline{b}$.
There are four individual terms.  Each term contains one of the
entries in the first factor (namely $a$ or $b$) and one of the entries
in the second factor (namely $\underline{a}$ or $\underline{b}$).  The
four terms represent thereby all possible ways of choosing one entry
in the first factor and one entry in the second factor.

\smallskip

\noindent 
ii) Multiply out the product $(a + b)(\underline{a}+ \underline{b})(A
+ B)$.  (Don't combine like terms yet.)  Does each term contain one
entry from the first factor, one from the second, and one from the
third?  How many terms did you get?  In fact there are two ways to
choose an entry from the first factor, two ways to choose an entry
from the second factor, and two ways to choose an entry from the third
factor.  Therefore, how many ways can you make a choice consisting of
one entry from the first, one from the second, and one from the third?
\smallskip

\noindent 
Now return to the general case:
	$$
(x + \Delta x)^n = \underbrace{(x + \Delta x)}_{1\mbox{-st}}
\underbrace{(x + \Delta x)}_{2\mbox{-nd}}
\underbrace{(x + \Delta x)}_{3\mbox{-rd}} \ldots \underbrace{(x + \Delta x)}_{n\mbox{-th}}
	$$
How many ways can you choose an entry from each factor and {\em not\/} get any $\Delta x$'s?  Multiply these chosen entries together; what does the product look like (apart from having no $\Delta x$'s in it)?

How many ways can you choose an entry from each factor in such a way
that the resulting product has {\em precisely one\/} $\Delta x$?
Describe all the various choices which give that result.  What does a
product that contains precisely one $\Delta x$ factor look like?  What
do you obtain for the sum of {\em all\/} such terms with precisely one
$\Delta x$ factor?

What is the minimum number of $\Delta x$ factors in any of the
remaining terms in the full expansion of $(x + \Delta x)^n$ ?

Do your calculations agree with this summary:  
	$$
(x + \Delta x)^n = x^n + nx^{n-1} \Delta x + R(\Delta x)^2 \ ?
	$$

\sub  Now find the value of  $\dfrac{f(x + \Delta x) - f(x)}{\Delta x}$.

\sub  Finally, find 
	$$
 \lim_{\Delta x \rightarrow 0} \dfrac{f(x + \Delta x) - f(x)}{\Delta x}.
	$$
Do you get $n x^{n-1}$?

\num In this problem we give another derivation of the power rule
based on writing
$$
x^{\textstyle r} = e^{\textstyle r \ln(x)}.
$$
Use the chain rule to differentiate $e^{\textstyle r \ln(x)}$.  Explain why your answer equals $rx^{r-1}$. 

\num Does the rule for the derivative of $x^r$ hold for $r=0$?  Why or
why not?

\num  In this exercise we prove the Addition Rule:  $F(x) = f(x) + g(x)$  implies  $F'(x) = f'(x) + g'(x)$.

\sub  Show  $F(x + \Delta x) - F(x) =  f(x + \Delta x) - f(x) + g(x + \Delta x) - g(x)$

\sub  Divide by  $\Delta x$  and finish the argument.

\num In this exercise we prove the Quotient Rule: $F(x) = f(x) / g(x)$
implies
	$$
	F'(x) = \frac{g(x)f'(x) - f(x)g'(x)}{(g(x))^2}
	$$

\sub Rewrite $F(x) = f(x) / g(x)$ as $f(x) = g(x) F(x)$.  Pretend for
the moment that you know what $F'(x)$ is and apply the Product Rule to
find $f'(x)$ in terms of $F(x)$, $g(x)$, $F'(x)$, $g'(x)$.  

\sub Replace $F(x)$ by $f(x) / g(x)$ in your expression for $f'(x)$ in
part a.

\sub Solve the equation in part b for $F'(x)$ in terms of $f(x)$,
$g(x)$, $f'(x)$ and $g'(x)$.

\num In this problem we calculate the derivative of $f(x) = x^n$ when
$n$ is a negative integer.  First write $n = -m$, so $m$ is a positive
integer.  Then $f(x) = x^{-m} = 1 / {x^m}$.

\sub Use the Quotient Rule and this new expression for $f$ to find
$f'(x)$ .

\sub  Do the algebra to re-express  $f'(x)$ as  $nx^{n-1}$.

\num In this problem we calculate the derivatives of $\sin{x}$ and
$\cos{x}$.  We will need the {\bf addition formulas}:
\begin{align*}
	\sin(A + B) &= \sin{A}\cos{B} + \cos{A}\sin{B} \\
	\cos(A + B) &= \cos{A}\cos{B} - \sin{A}\sin{B}
\end{align*}

\noindent First tackle  $f(x) = \sin{x}$:

\sub Use the addition formula for $\sin(A + B)$ to rewrite $f(x +
\Delta x)$ in terms of $\sin(x)$, $\cos(x)$, $\sin(\Delta x)$, and
$\cos(\Delta x)$.  \sub The quotient $\dfrac{f(x + \Delta x) -
  f(x)}{\Delta x}$ can now be written in the form
	$$
	P(\Delta x) \cdot \sin{x} + Q(\Delta x) \cdot \cos{x} \, ,
	$$
where $P$ and $Q$ are specific functions of $\Delta x$.  What are the formulas for those functions?

\sub  Use a calculator or computer to estimate the limits
	$$
	\lim_{\Delta x \to 0}P(\Delta x) \qquad \mbox{and} \qquad
		\lim_{\Delta x \to 0}Q(\Delta x) \ .
	$$
(Try  $\Delta x =$ .1, .01, .001, .0001.  Be sure your calculator is set on radians, not degrees.)  Using part b you should now be able to determine the limit
	$$
	\lim_{\Delta x \rightarrow 0} 
		\frac{f(x + \Delta x) - f(x)}{\Delta x} \ 
	$$  
by writing it in the form
	$$
	\left( \lim_{\Delta x \rightarrow 0} P(\Delta x) \right) 
					\cdot \sin{x} + 
	\left( \lim_{\Delta x \rightarrow 0} Q(\Delta x) \right) 
					\cdot \cos{x} \ .
	$$

\sub  What is  $f'(x)$?

\sub  Proceed similarly to find  the derivative of $g(x) = \cos{x}$.


\num In this problem we calculate the derivatives of the other
circular functions.  Use the quotient rule together with the
derivatives of $\sin{x}$ and $\cos{x}$ to verify that the derivatives
of the other four circular functions are as given in the table below:
$$
\begin{array}{ccc}
\mbox{function}  & & \mbox{derivative} \\ [2pt] 
\hline \\ [-9pt]
\tan{x} = \dfrac{\sin{x}}{\cos{x}}      
		& &       \sec^2{x} \\ [.15in]
\csc{x} = \dfrac{1}{\sin{x}}     
		 & & -\cot{x}\csc{x}  \\ [.15in]
\sec{x} = \dfrac{1}{\cos{x}}     
		&   &    \sec{x}\tan{x} \\ [.15in]
\cot{x} = \dfrac{1}{\tan{x}}	 
		&    &     -\csc^2{x}
\end{array}
$$

\subsubsection{Differential Equations}

\vspace{-10pt}

\num If $y = f(x)$ then the {\bf second derivative}\index{second
  derivative}\index{derivate!second} of $f$ is just the derivative of the derivative of $f$;
it is denoted $f''(x)$ or $d^2y/dx^2$.  Find the second derivative of
each of the following functions.

\sub $f(x) = e^{3x-2}$

\sub $f(x) = \sin{\omega x}$, where $\omega$ is a constant

\sub $f(x) = x^2e^x$

\num Show that $e^{3x}$ and $e^{-3x}$ both satisfy the ({\em second
  order\/}) differential equation
	$$
	f''(x) = 9 f(x).
	$$
Furthermore, show that {\em any\/} function of the form $g(x) = \alpha e^{3x} + \beta e^{-3x}$ satisfies this differential equation.  Here $\alpha$ and $\beta$ are arbitrary constants.  Finally, choose $\alpha$ and $\beta$ so that $g(x)$ also satisfies the two conditions $g(0) = 12$ and $g'(0) = 15$.

\num Show that $y = \sin{x}$ satisfies the differential equation $y''
+ y = 0$.  Show that $y = \cos{x}$ also satisfies the differential
equation.  Show that, in fact, $y=a \sin{x} + b \cos{x}$ satisfies the
differential equation for any choice of constants $a$ and $b$.  Can
you find a function $g(x)$ that satisfies these three conditions:
\begin{align*}
	g''(x) + g(x) & =  0 \\
		g(0)  & =  1 \\
		g'(0) & =  4 ?
\end{align*}

\num Show that $\sin{\omega x}$ satisfies the differential equation
$y'' + {\omega}^2 y = 0$.  What other solutions can you find to this
differential equation?  Can you find a function $L(x)$ that satisfies
these three conditions:
\begin{align*}
L''(x) + 4L(x) & =  0 \\
	L(0)   & =  36 \\
	L'(0)  & =  64 ?
\end{align*}


\subsubsection{The Colorado River Problem}.
  
\vspace{-10pt}

Make your answer to this sequence of questions an essay.  Identify all
the variables you consider (e.g., ``$A$ stands for the area of the
lake''), and indicate the functional relationships between them (``$A$
depends on time $t$, measured in weeks from the present'').  Identify
the derivatives of those functions, as necessary.

The Colorado River---which excavated the Grand Canyon, among others---used to empty into the Gulf of California.  It no longer does.  Instead, it runs into a marshy area some miles from the Gulf and stops.  One of the major reasons for this change is the construction of dams---notably the Hoover Dam.  Every dam creates a lake behind it, and every lake increases the total surface area of the river.  Since the rate at which water evaporates is proportional to the area of the water surface exposed to air, the lakes along the Colorado have increased the loss of river water through evaporation.  Over the years, these losses (in conjunction with other factors, like increased usage by a rapidly growing population) have been significant enough to dry up the river at its mouth. 
\medskip

\noindent
\begin{minipage}{3.5in}
\num Let us analyze the evaporation rate along a river that was
recently dammed.  Suppose the lake is currently 50 yards wide, and
getting wider at a rate of 3 yards per week.  As the lake fills, it
gets longer, too.  Suppose it is currently 950 yards long, and it is
extending upstream at a rate of 15 yards per week. Assuming the lake
remains approximately rectangular as it grows, find
\end{minipage} \
\begin{minipage}{2in}
\vspace*{1.75in}

\special{psfile=../ps/calc1/lake.ps}
\end{minipage}

\noindent 
\sub  the current area of the lake, in square yards;

\sub the rate at which the surface of the lake is currently growing,
in square yards per week.

\num Suppose the lake continues to spread sideways at the rate of 3
yards per week, and it continues to extend upstream at the rate of 15
yards per week. \sub Express the area of the lake as a (quadratic!)
function of time, where time is measured from the present, in weeks,
and where the lake's area is as given in problem 25.

\sub How many weeks will it take for the lake to cover 30 acres (=
145,200 square yards)?

\sub  At what rate is the lake surface growing when it covers 30 acres?


\num Compare the rates at which the surface of the lake is growing in
problem 25 (which is the ``current'' rate) and in problem 26 (which is
the rate when the lake covers 30 acres).  Are these rates the same?
If they are not, how do you account for the difference?  In
particular, the width and length grow at fixed rates, so why doesn't
the area?  Use what you know about derivatives to answer the question.

\num Suppose the local climate causes water to evaporate from the
surface of the lake at the rate of 0.22 cubic yards per week, for each
square yard of surface.  Write a formula that expresses total
evaporation per week in terms of area.  Use $E$ to denote total
evaporation.

\num The lake is fed by the river, and that in turn is fed by
rainwater and groundwater from its watershed.  (The {\bf watershed},
or basin, of a river is that part of the countryside containing the
ponds and streams which drain into the river.)  Suppose the watershed
provides the lake, on average, with 25,000 cubic yards of new water
each week.

Assuming, as we did in problem 25, that the lake widens at the
constant rate of 3 yards per week, and lengthens at the rate of 15
yards per week, will the time ever come that the water being added to
the lake from its watershed balances the water being removed by
evaporation?  In other words, will the lake ever stop filling?

\newpage
%%%%%%%%%%%%%%%%%%%%%%%%%%%%%%%%%%%%%%%%%%%%%%%%%%%%%%%%%%%%%%%%%%%%%%%%%
%                                                                       %
%                                Section 2                              %
%                                                                       %
%%%%%%%%%%%%%%%%%%%%%%%%%%%%%%%%%%%%%%%%%%%%%%%%%%%%%%%%%%%%%%%%%%%%%%%%%     

\section{Finding Partial Derivatives} 

     We know from  Chapter 3 that no additional formulas are needed to calculate {\em partial\/} derivatives.  We simply use the usual differentiation formulas, treating all the variables except one---the one with respect to which the partial derivative is formed---as if they were constants.  If we do this we get new techniques for analyzing rates of change in problems that involve functions of several variables.

\subsection{Some Examples}

      Here are two examples to illustrate the technique for calculating partial derivatives:
      \mar{\vspace{20pt}Finding formulas for partial derivatives}

\begin{enumerate}

\item Suppose $f(x,y) = x^2y + 5x^3 - \sqrt{x+y}$.  Then
\begin{align*}
f_x(x,y) &= 2xy + 15x^2 - \frac{1}{2\sqrt{x+y}}, \quad\mbox{and}\\
f_y(x,y) &= x^2  - \frac{1}{2\sqrt{x+y}}.
\end{align*}

\item Suppose $g(u,v) = e^{\displaystyle uv} + \displaystyle{\frac{u}{v}}$.  Then
\begin{align*}
g_u(u,v) &= ve^{\displaystyle uv} + \frac{1}{v}, \quad\mbox{and}\\[8pt]
g_v(u,v) &= ue^{\displaystyle uv} - \frac{u}{v^2}.
\end{align*}

\end{enumerate}

\subsection{Eradication of Disease}

Controlling---or, better still, eradicating---a communicable disease
depends first on the development of a vaccine.  But even after this
step has been accomplished, public health officials must still answer
important questions, including:

\begin{itemize}

\item What proportion of the population must be vaccinated in order to
  eliminate the disease?

\item  At what age should people be vaccinated? 

\end{itemize}

In their 1982 article, ``Directly Transmitted Infectious Diseases:
Control by Vaccination,'' ({\em Science}, Vol. 215, 1053--1060), Roy
Anderson\index{Anderson, R.} and Robert May\index{May, R.} formulate a model
for the spread of disease that permits them to answer these and other
questions.  For a particular disease in a particular environment, the
important variables in their model are
\begin{enumerate}

\item  The average human life expectancy $L$, in years;

\item The average age $A$ at which individuals catch the disease, in
  years;

\item The average age $V$ at which individuals are vaccinated against
  the disease, in years.
\end{enumerate}

Anderson and May deduce from their model that in order to eradicate
the disease, the proportion of the population that is vaccinated must
exceed $p$, where $p$ is given by
$$
p = \frac{L + V}{L + A}.
$$ 

For a disease like measles, public health officials can directly
affect the variable $V$,
%
\mar{Partial derivatives can tell us which variables are most significant}     
%
for example by the recommendations they make to physicians about
immunization schedules for children.  They may also indirectly affect
the variables $A$ and $L$, because public health policy influences
factors which can modify the age at which children catch the disease
or the overall life expectancy of the population.  (Many other factors
affect these variables as well.)  Which of these three variables has
the greatest effect on the proportion of the population that must be
vaccinated?
	
In other words, which is largest: $\partial p/\partial L$, $\partial
p/\partial A$, or $\partial p/\partial V$?

     Using the rules, we compute:
\begin{align*}
\frac{\partial p}{\partial L} &= 
   \frac{1\cdot (L+A)-1\cdot (L+V)}{(L+A)^2} = \frac{A-V}{(L+A)^2},\\
\frac{\partial p}{\partial A} &= \frac{-(L+V)}{(L+A)^2}, \quad\mbox{and}\\
\frac{\partial p}{\partial V} &= \frac{1}{L+A}. 
\end{align*}
For measles in the United States, reasonable values of the variables are $L$ = 70 years, $A$ = 5 years and $V$ = 1 year.  Using these values, the crucial proportion of the population needing to be vaccinated is $p=71/75 = .947$, and the partial derivatives are
\begin{alignat*}{2}
\dfrac{\partial p}{\partial L} &= 
    \dfrac{4}{(75)^2}   &&= .0007, \\[4pt]
\dfrac{\partial p}{\partial A} &= 
    \dfrac{-71}{(75)^2} &&= -.0126,   \\[4pt]
\dfrac{\partial p}{\partial V} &= 
    \dfrac{1}{75} &&= .0133. 
\end{alignat*}


\bigskip
A comment is in order here on units.  
	\mar{Determining units}
While the input variables $L\/$, $A\/$ and $V$ are all measured in years---so the rates are {\em per year\/}, the output variable $p$ is dimensionless: it is the  ratio of persons vaccinated to persons not vaccinated.  It would be reasonable to write $p$ as a percentage.  Then we can attach the units {\em percent per year\/} to each of the three partial derivatives.  Thus we have:
\begin{alignat*}{2}
\frac{\partial p}{\partial L} 
          &= .07\% \hspace{.25in} &\mbox{per year} \\[4pt]
\frac{\partial p}{\partial A} 
          &= -1.26\% \hspace{.25in} &\mbox{per year} \\[4pt]
\frac{\partial p}{\partial V} 
          &= 1.33\%  \hspace{.25in} &\mbox{per year.}
\end{alignat*}

    It is not surprising that a change in average life 
expectancy has a negligible effect on the proportion $p$ of 
the population that must be vaccinated in order to eradicate 
measles.  Nor is it surprising that changing the age of 
vaccination has the greatest effect on $p$.  But it is not 
obvious ahead of time that changing the age at which 
children catch the disease has nearly as large an effect on 
$p$:
\begin{itemize}
\item  Decreasing the age of vaccination decreases the proportion $p$ by 1.33\% per year of decrease.
\item  Increasing the age at which children catch measles decreases the proportion $p$ by 1.26\% per year of increase.
\end{itemize}

Changes can also go the ``wrong'' way.  For example, in an area where
use of communal child care facilities is growing, contact among very
young children increases, and the age at which children are exposed
to---and can catch---communicable diseases like measles falls.  The
Anderson--May model tells us that immunization practices must change to
compensate: either the age of vaccination must drop a like amount, or
the fraction of the population that is vaccinated must grow by 1.26\%
per year of decrease in the average age at infection.

\subsection{Exercises}

\subsubsection{Finding Partial Derivatives}

\vspace{-10pt}
\num  Find the partial derivatives of the following functions.
\sub $x^2 y$.
\sub $\sqrt{x + y}$
\sub $e^{xy}$
\sub $\displaystyle{\frac{y}{x}}$
\sub $\dfrac{x + y}{y + z}$
\sub $\sin{\displaystyle{\frac{y}{x}}}$

\numa Suppose $f(x,y) = \displaystyle{e^{-(x+2y)}(2x - 5y)}$.  Find
$f_x(x,y)$ and $f_y(x,y)$.

\sub Find a point $(a, b)$ at which $f_x(a,b)=0$.  At such a point a
small change in $x$ leaves the value of $f$ virtually unchanged.

\sub Find a point $(a, b)$ at which a small {\em increase} in the
$x$-value would produce the same change in $f(a,b)$ as would the
same-sized {\em decrease} in the $y$-value.

\num Suppose $g(u,v) = \dfrac{\sin{u+v^2}+7uv}{1+u^2+v^4}$.  Find
$g_u(u,v)$ and $g_v(u,v)$.

\num The {\bf second partial derivatives}\index{derivative!second partial} of $z
= f(x, y)$ are the partial derivatives of $\partial f/\partial x$ and
$\partial f/\partial y$, namely:
\begin{align*}
\dfrac{\partial ^2 f}{\partial x ^2} & = 
		\dfrac{\partial}{\partial x} \left( 
		\dfrac{\partial f}{\partial x} \right) \\[4pt]
\dfrac{\partial ^2 f}{\partial x \partial y} & = 
		\dfrac{\partial}{\partial x} \left( 
		\dfrac{\partial f}{\partial y} \right) \\[4pt]
\dfrac{\partial ^2 f}{\partial y ^2} & = 
		\dfrac{\partial}{\partial y} \left( 
		\dfrac{\partial f}{\partial y} \right) 
\end{align*}
Find the three second partial derivatives of the following functions.
\sub $x^2 y$.
\sub $\sqrt{x + y}$
\sub $e^{xy}$
\sub $\displaystyle{\frac{y}{x}}$
\sub $\sin{\displaystyle{\frac{y}{x}}}$

\subsubsection{Eradication of Disease}

\vspace{-10pt}
\num Suppose you were dealing with measles in a developing country where $L = 50$ years, $A = 4$ years, and $L=2$ years.  Discuss the impact on measles control if increased public health efforts increase $L$ to  55 years, $A$ to 5 years, and decrease $V$ to 1.5 years.

\subsubsection{Partial differential equations}\index{partial differential equations}\index{differential equation!partial}

\vspace{-10pt}
\num  Show that the function 
$z = \dfrac{1}{\sqrt{t}}
\exp{\dfrac{-x^2}{4t}}$ satisfies the {\bf partial differential equation\/}
	$$
	\dfrac{\partial ^2 z}{\partial x^2} = 
		\dfrac{\partial z}{\partial t}.
	$$

\num  Show that {\em every\/} linear function of the form
$z = px + qy + c$ satisfies the partial differential equation
	$$
	\dfrac{\partial ^2 z}{\partial x^2} +
		\dfrac{\partial ^2 z}{\partial y^2} = 0.
	$$
Here $p$, $q$, and $c$ are arbitrary constants.

\num  Show that the function $z = e^x\sin{y}$ also satisfies the partial differential equation
	$$
	\dfrac{\partial ^2 z}{\partial x^2} +
		\dfrac{\partial ^2 z}{\partial y^2} = 0.
	$$

%%%%%%%%%%%%%%%%%%%%%%%%%%%%%%%%%%%%%%%%%%%%%%%%%%%%%%%%%%%%%%%%%%%%%%%%%
%                                                                       %
%                              Section 3                                %
%                                                                       %
%%%%%%%%%%%%%%%%%%%%%%%%%%%%%%%%%%%%%%%%%%%%%%%%%%%%%%%%%%%%%%%%%%%%%%%%%     


\section{The Shape of the Graph of a \mbox{Function}}

We know from chapter 3 that the derivative gives us qualitative
information about the shape of the graph of a differentiable function.

\vspace{-4pt}
\begin{center}
\begin{tabular}{ll} 
function &  derivative \\ \hline
increasing  &  positive \\
decreasing  & negative \\
level &  zero \\
steep (rising or falling) & large (positive or negative) \\
gradual (rising or falling) \qquad & small (positive or negative) \\
straight & constant
\end{tabular}
\end{center}
\vspace{-4pt}

Having a formula for the derivative of a function will thus give us a
great deal of information about the behavior of the function itself.
In particular we will be interested in using the derivative to solve
%
\mar{Contexts for optimization problems}	
%
{\bf optimization problems}\index{optimization}---finding maximum or
minimum values of a function.  Such problems occur frequently in many
fields.
	
\begin{itemize}
\item Economists actually define human rationality in terms of
  optimization.  Each person is assumed to have a {\em utility
    function}, a function that assigns to each of many possible
  outcomes its {\em utility}, a numerical measure of its value to her.
  (Different people may have different utility functions, depending on
  their personal value systems.)  A rational person is one who acts to
  maximize her utility.  Some utilities are expressed in terms of
  money.  For example, a rational manufacturer will seek to maximize
  her profit (in dollars).  Her profit will depend on---that is, be a
  function of---such variables as the cost of her raw materials and
  the unit price she charges for her product.

\item Many physical laws are expressed as minimum principles.
  Ordinary soap bubbles exhibit one of these principles.  A soap film
  has a {\em surface energy\/} which is proportional to its surface
  area.  For almost any physical system, its stable state is one which
  minimizes its energy.  Stable soap films are thus examples of {\em
    minimal surfaces}.  Interfaces involving crystals also have
  surface energies, leading to the study of crystalline minimal
  surfaces.

\item Statisticians develop mathematical summaries for data---in other
  words, mathematical models.  For example, a relationship between two
  numerical variables may be summarized by a linear function, say
  $y=mx+b$, where $x$ and $y$ are the variables of interest.  It would
  be very rare to find data that were exactly linear.  In a particular
  case, the statistician chooses the linear model that minimizes the
  {\em discrepancy\/} between the actual values of $y$ and the
  theoretical values obtained from the linear function.  Statisticians
  frequently measure this discrepancy by summing the squares of the
  differences between the actual and the theoretical values of $y$ for
  each data point.  The {\em best-fitting line\/} or {\em regression
    line\/} is the graph of the linear function which is optimal in
  this sense.

\item Psychologists who study decision-making have found that some
  people are ``risk-averse"; they make their decisions primarily to
  avoid risks.  If we regard risk as a function of the various
  outcomes under consideration (a bit like a utility function), such a
  person acts to minimize this function.

\end{itemize}	
	
		
The derivative is the key tool here. We will develop a general
procedure for using the derivative of a function to locate its
extremes.

\subsection{Language}

Here is a graph of what we might consider a ``generic" differentiable
function.\label{`concave'}

\vspace{130pt}

\special{psfile=../ps/calc1/pic244a.ps}

\noindent
The most distinctive features are the hill tops and valley bottoms,
points where the graph levels and the derivative is zero.  \mar{Local
  extremes and global extremes} We distinguish between {\bf local}
extremes\index{extreme!local}, like those occurring at the points
$x=b$, $x=c$, and $x=d$ and a {\bf global}
extreme\index{extreme!global}, like the global maximum at the point
$x=a$.  The function has a {\bf local} {\bf minimum} at $x=b$ because
$f(x) \geq f(b)$ for all $x$ {\em sufficiently near\/} $b$.  The
function has a {\bf global} {\bf maximum} at $x=a$ because $f(x) \leq
f(a)$ for {\em all\/} $x$.  Notice that this particular function does
not have a global minimum. What kinds of local extremes does the
function have at $x=c$ and $x=d$?  The convention is to say that {\em
  all\/} extremes are local extremes, and a local extreme may or may
not also be a global extreme.


Examining the graph of as simple a function 
%
\mar{A function may have no extreme at a point where
  its derivative equals zero}
%
as $f(x)=x^3$ shows us that a function need not have any extremes at
all.  Moreover, since $f'(0)=0$ for this function, a zero derivative
doesn't necessarily identify a point where an extreme occurs.

\vspace{110pt}

\special{psfile=../ps/calc1/pic244b.ps}


Can a function have an extreme at a point 
%
\mar{An extreme can\\ occur at a cusp}
%
{\em other\/} than where the derivative is zero?  Consider the graph
of $f(x)=x^{2/3}$ below.

\vspace{90pt}

\special{psfile=../ps/calc1/pic245.ps} 

\noindent 
This function is differentiable everywhere except at the point $x=0$.
And it is at this very point, where
$$
f'(x) = \frac{2}{3}x^{-1/3} = \frac{2}{3 \, \sqrt[3]{x}}
$$ 
is undefined, that the function has its global minimum.  For this
reason, points where the derivative fails to exist (or is infinite)
are as important as points where the derivative equals zero.  All of
these kinds of points are called {\bf critical points}\index{critical
  point}\label{crit-begin} for the function.
\begin{center}
\begin{picture}(310,45)
\put(-10,0){\framebox(330,45)
{\shortstack {A {\bf critical point} for a function $f$ is a point on the\\
graph of $f$ where $f'$ equals zero or infinity or fails to exist.}}}
\end{picture}
\end{center}

How can we recognize, from looking at its graph, that a function fails
to be differentiable at a point?  We know that when the graph has a
sharp corner or {\bf cusp}\index{cusp} it isn't locally linear at that
point, and so has no derivative there. Thus critical points occur at
these places.  When the graph is locally linear but vertical, the
slope cannot be given a finite numerical value. Critical points also
occur at these vertical places where the slope is infinite.  The graph
below is an example of such a curve.


\vspace{110pt}

\special{psfile=../ps/calc1/pic246.ps}

\label{continuous-begin}
A weaker condition than differentiability, 
%
\mar{Continuous functions}
%
but one that is useful, especially in this context, is {\bf
  continuity}.\index{continuity} We say that a function $f$ is {\bf
  continuous at a point}\index{function!continuous} $x=a$ if
 \begin{itemize}

\item it is defined at the point, and 

\item we can achieve changes in the output that are arbitrarily small
  by restricting changes in the input to be sufficiently small.

\end{itemize}

This second condition can also be expressed in the following form:
\begin{quote}
Given any positive number $\epsilon$ (the proposed limit on the change
in the output is traditionally designated by the Greek letter
$\epsilon$, pronounced `epsilon'), there is always a positive number
$\delta$ (the Greek letter `delta'), such that whenever the change in
the input is less than $\delta$, then the corresponding change in the
output will be less than $\epsilon$.
\end{quote} 

A function is continuous on a set of real numbers if it is continuous
at each point of the set.
%
\mar{Graphs of continuous functions have\\ no gaps or jumps}
%
A natural way to think of (and recognize) a function which is
continuous on an interval is by its graph on that interval, which is
continuous in the usual sense: it has no gaps or jumps in it.  You can
draw it without picking up your pencil.  Of the four functions whose
graphs appear on the next page, $f$ is continuous (and
differentiable); $g$ is continuous, but not differentiable (because of
the cusp at $x=a$); $h$ is not continuous, because $h$ is undefined at
$x=b$; and $k$ is not continuous, because of the ``jump" at $x=c$.

	Luckily, the functions that we are likely to encounter are continuous on their natural domains.  
	\mar{A quotient isn't continuous where its denominator is zero}
Among functions given by formulas, the only exceptions we have to worry about are quotients, like $f(x)=1/x\/$, which have gaps in their natural domains at points where the denominator vanishes (at $x=0$ in the case of $f(x)=1/x\/$), so their domains are not single intervals.  For convenience, we usually confine our attention to functions continuous on an interval.

\vspace{120pt}

\special{psfile=../ps/calc1/pic247a.ps}

\vspace{130pt}

\special{psfile=../ps/calc1/pic247b.ps}



\subsection{The Existence of Extremes}

Not every function has extremes, as the example of $f(x)=x^3$ shows.
We are, however, guaranteed\index{existence of extremes} extremes for
certain functions.

\medskip
\noindent
{\bf Principle I.}
We are guaranteed that a function 
	\mar{A continuous function on a finite closed interval has extremes}
has a global maximum and a global minimum if we know:
\begin{itemize}

\item the domain of the function is a {\em finite closed interval},
  and

\item the function is {\em continuous\/} on this interval. 
\end{itemize}
The domain restriction in most optimization problems is likely to come
from the physical constraints of the problem, not the mathematics.  A
{\bf finite closed interval}\index{interval!finite closed}, written
$[a,b]$, is the set of all real numbers between $a$ and $b$, including
the endpoints $a$ and $b$.  Other kinds of finite intervals are
$(a,b]$, (this excludes the endpoint $a$), $[a,b)$, and $(a,b)$.
    Infinite intervals include open and closed ``rays" like
    $(a,+\infty)$, $[a, +\infty)$, $(-\infty, b)$ and $(-\infty,b]$,
    and the entire real line.

To illustrate Principle~I, we note first that it does not apply to
$f(x)=1/x$ on the finite closed interval $[-1,2]$, because this
function isn't continuous at every point on this interval---in fact,
the function isn't even defined for $x=0$.

\vspace{125pt}

\special{psfile=../ps/calc1/pic248a.ps}

\special{psfile=../ps/calc1/pic248b.ps}

\noindent
By contrast, Principle I {\em does\/} apply to the same function if we
change its domain to a finite closed interval that does not include $x=0$.
In the figure above we use $[1,3]$.


A function hat fails to satisfy the first condition of Principle~I
can still have global extremes.  For any function continuous on an
interval we can apply the following principle.
	
\medskip
\noindent
{\bf Principle II.}  If a function $f$ is continuous 
	\mar{Continuous functions have extremes\\ only at endpoints\\ or critical	points}
on an interval and has a local extreme at a point $x=c$ of the interval, then either $x=c$ is a critical point for the function, or $x=c$ is an endpoint of the interval.  In other words, we are guaranteed that one of the following three conditions holds:
\begin{itemize}
\item $f'(c)=0$, 
\item $f'(c)$ is undefined,
\item $x=c$ is an endpoint of the interval.
\end{itemize}

\newpage

	The following graphs illustrate three local maxima, each satisfying one of these three conditions.
	
\vspace{80pt}

\special{psfile=../ps/calc1/pic249.ps}

When we apply Principle II to optimization problems, an important part
of the task will be ascertaining which, if any, of the critical points
or endpoints we find actually gives the extreme we're looking for.
We'll examine a variety of techniques, graphical and analytical, for
locating critical points and determining what kind of extreme point
(if any) they are.
\label{continuous-end}	
	

\subsection{Finding Extremes}

\subsubsection{Using a graphical approach}

	If we can use a computer 
	\mar{Computer graphing can be easy if the general location of the extremes is known}
to examine the graph of the function of interest, we can determine the existence and location of extremes by inspection.  However, every graphing utility requires the user to specify the interval on which the function will be graphed, and careful analysis may be required in order to choose an interval that contains all the extremes of interest.

For functions given by data, whose graphs have only finitely many points, we can zoom in to find the exact coordinates of the extreme datapoints.  For a function given by a formula, we can estimate the coordinates of an extreme to arbitrary accuracy by zooming in on the point as closely as desired.  This is the method we used in some of the exercises of Chapter 1, and it is quite satisfactory in many situations.  

\subsubsection{Using the formula for the derivative.}

\enlargethispage*{20pt}
	In this chapter we are concentrating 
	\mar{Formulas can give exact answers and can handle parameters}
on functions given by formulas.  For these functions we may want a method other than the approximation using a graphing utility described above. 
	\begin{itemize}
\item For some functions, the determination of extremes using a formula for the derivative is at least as easy as using the computer.
\item  Some functions are described in terms of a {\em parameter\/}, a constant whose value may vary from one problem to another.  For example, the rate equation from the S-I-R model for change in the number of infected, $I' = aSI-bI\/$, involves two parameters $a$ and $b\/$, the transmission and recovery coefficients.  For such a function, we cannot use the computer unless we specify a numerical value for each parameter, thus limiting the generality of our results.
\item  The computer gives only an approximation, while a precise answer may convey important additional information.
\end{itemize}

	We will assume that the function we are studying is continuous on the domain of interest and that its domain is an interval.   Our procedure for finding extremes of a function given by a formula $y=f(x)$ is thus a direct application of \mbox{Principle II.}   

\subsubsection{The search for local extremes}

\begin{enumerate}
\item Determine the domain of the function and identify its endpoints, if any.  Keep in mind that in an applied context, the domain may be determined by physical or other restrictions.  
\item  Find a formula for $f'(x)$.
\item  Find any roots of $f'(x) = 0$ in the domain.  An estimation procedure, for example Newton's Method (see the end of this chapter),  may be used for complicated equations.
\item  Find any points in the domain where $f'(x)$ is undefined.
\item  Determine the shape of the graph to locate any local extremes.  Find the shape either by looking directly at the graph of $y = f(x)$ or by analyzing the sign of $f'(x)$ on either side of each critical point $x=c$.  (Sometimes the second derivative $f''(c)$ aids the analysis; see the exercises for details.)
\end{enumerate}

\noindent {\bf The search for global extremes}
\begin{enumerate}
\item  Find the local extremes, as above.
\item  If the domain is a finite closed interval, it is only necessary to compare the values of $f$ at each of the critical points and at the endpoints to determine which is the global maximum and which is the global minimum.
\item  More generally, use the shape of the graph to ascertain whether the desired global extreme exists and to identify it.
\end{enumerate}

In the succeeding sections and in the exercises we will carry out
these search procedures in a variety of situations. Since there are
often substantial algebraic difficulties in analyzing the sign of the
derivative, or even determining when it is equal to zero, in realistic
problems, in section~5 we will develop some numerical methods for
handling such complications.
\label{crit-end}


\subsection{Exercises}

\subsubsection{Describing functions}

\vspace{-10pt}
\num  For each of the following graphs, is the function continuous?
 Is the function differentiable?

\vspace{45pt}
 
\special{psfile=../ps/calc1/pic252a.ps}
\special{psfile=../ps/calc1/pic252c.ps}
\special{psfile=../ps/calc1/pic252b.ps}
\special{psfile=../ps/calc1/pic252d.ps}

\num For each of the following graphs of a function $y=f(x)$, is $f'$
increasing or decreasing?  At the indicated point, what is the sign of
$f'$?  Is $f'$ ({\em not} $f$) increasing or decreasing at the point?
What does this then say about the sign of $f''$ (the derivative of
$f'$) at the point?

\vspace{110pt}
\special{psfile=../ps/calc1/pic253a.ps}

\vspace{110pt}
\special{psfile=../ps/calc1/pic253c.ps}

\num For each of the following, sketch a graph of $y=f(x)$ that is
consistent with the given information.  On each graph, mark any
critical points or extremes.

\sub $f'(x) > 0$ for $x<1\/$; $f'(1)=0\/$; $f'(x) < 0$ for $1<x<2\/$;
$f'(2)=0\/$; $f'(x)>0$ for $x>2\/$.

\sub  $f'(x) > 0$ for $x<2\/$; $f'(2)=0\/$; $f'(x)>0$ for $x>2\/$.

\sub  $f'(x) > 0$ for $x<2\/$; $f'(2)=0\/$; $f'(x)<0$ for $x>2\/$.

\sub  $f'(3)=0\/$; $f''(3) > 0\/$.

\num  {\bf The geometric meaning of the second derivative}  If $f$ is any function and if $f''(x)$ is positive over some interval, then $f'$ is increasing over that interval, and we say the curve is {\bf concave upward}\index{concave upward} over that interval.  If $f''(x)$ is negative over some interval, then $f'$ is decreasing over that interval, and we say the curve is {\bf concave downward}\index{concave downward} over that interval.  You should study the graphs in problem 2 until you are clear what this means geometrically.  
\sub Suppose there were four functions $f$, $g$, $h$, and $k$, and that $f(0)=g(0)=h(0)=k(0)=0$, and $f'(0)=g'(0)=h'(0)=k'(0)=1$.  Suppose, moreover, that $f''(0)=1$, $g''(0)=5$, $h''(0)=-1$, and $k''(0)=-5$.  Sketch possible graphs of these functions near the origin.
\sub We know that the magnitude of the first derivative tells us how steep the curve is---the greater the value of $f'$, positive or negative, the steeper the curve. What does the magnitude of the second derivative tell us geometrically about the shape of the curve?  Complete the sentence: ``The greater the value of $f''$, the \rule[-2pt]{1.25in}{.2pt}.''

\numa Here is the graph you saw back on page~\pageref{`concave'}:

\vspace{130pt}

\special{psfile=../ps/calc1/pic244a.ps}

\noindent
At which point is the second derivative greater, $b$ or $d$, and why?
\sub Reproduce a sketch of this curve and indicate where the curve is concave up and where it is concave down.
\sub What must be true about the second derivative at the points where the curve changes concavity from up to down, or vice versa?  Give a clear justification for your answer.
\sub  What must be true about the second derivative near the right-hand end of the graph, and why?
\sub Put all this together to sketch the graph of the second derivative of this function.  Label the values $a$ -- $d$ on your sketch.

\num {\bf Second derivative test for maxima and minima}\index{second derivative test}
Explain why the following test works.
\begin{itemize}
\item  If $f'(c)=0$ and $f''(c) > 0\/$, then $f$ has a {\em local minimum\/} at $x=c\/$.  

\item If $f'(c)=0$ and $f''(c)<0\/$, then $f$ has a {\em local maximum\/} at $x=c\/$.
\end{itemize}
(Hint: If $f'(c)=0$ and $f''(c) > 0\/$, what can you say about the value of $f'(x)$---and hence the geometry of the graph of $f$---on either side of $c$?)

\subsubsection{Finding critical points}

\vspace{-10pt}
\num  For each of the following functions, find the critical points, if any, without using a computer or calculator.  Can you sketch the graph of the function near the critical point?  Use the second derivative test if you can't figure the behavior out from a simpler inspection.
\sub  $f(x) = x^{1/3}$
\sub  $f(x) = x^3 + \displaystyle{\frac{3}{2}}x^2 - 6x + 5$
\sub  $f(x) = \displaystyle{\frac{2x+1}{x-1}}$
\sub  $f(x)=\sin{x}$
\sub  $f(x) = \sqrt{1-x^2}$
\sub $f(x) = \displaystyle{\frac{e^x}{x}}$
\sub $f(x) = x\ln x$
\sub $f(x)=\displaystyle{x^c +\frac{1}{x^c}}$ where c is some constant 
  

\num  For each of the following graphs, mark any critical points or extremes.  Indicate which extremes are local and which are global.  (Assume that at their ends the curves continue in the direction they are headed.)

\vspace*{1.5in} 
 \special{psfile=../ps/calc1/pic254a.ps}

 \vspace*{1.5in}
  \special{psfile=../ps/calc1/pic254b.ps}

\vspace*{1.5in}
  \special{psfile=../ps/calc1/pic254c.ps}
 
\subsubsection{Finding extremes}

\vspace{-5pt} 

Except where indicated, you should not use a computer or calculator to
solve the following problems.

\num For what positive value of $x$ does $f(x) = x + \dfrac{7}{x}$
attain its minimum value?  Explain how you found this value.

\num For what value of $x$ in the interval $[1,2]$ does $f(x) = x +
\dfrac{7}{x}$ attain its minimum value?  Explain how you found this
value.

\num Use a graphing program to make a sketch of the function $y = f(x)
= x^{2} 2^{-x}$ on the interval $0 \leq x \leq 10$.  From the graph,
estimate the value of $x$ which makes $y$ largest, accurately to four
decimal places.  Then find where $y$ takes on its maximum by setting
the derivative $f'(x)$ equal to 0.  You should find that the maximum
occurs at $x = 2/\ln{2}$.  Finally, what is the numerical value of
this estimate of $2/\ln{2}$, accurate to seven decimal places?

\num Use a graphing program to sketch the graph of $y = \dfrac{1}{x^2}
+ x$ on the interval $0 < x < 4$ and estimate the value of $x$ that
makes $y$ smallest on this interval.  Then use the derivative of $y$
to find the exact value of $x$ that makes $y$ smallest on the same
interval.  Compare the estimated and exact values.

\num What is the smallest value $y = \dfrac{4}{x^2} + x$ takes on when
$x$ is a positive number?  Explain how you found this value.

\num The function $y = x^4 - 42 x^2 -80x$ has two local minima.  Where
are they (that is, what are their $x$ coordinates), and what are their
values?  Which is the lower of the two minima?  In this problem you
will have to solve a cubic equation.  You can use an estimation
procedure, but it is also possible to solve by factoring the cubic.
(The roots are integers.)

\num The function $y = x^4 - 6x^2 + 7$ has two local minima.  Where
are they, and which of the two is lower?  In this problem you will
have to solve a cubic equation.  Do this by factoring and then
approximating the $x$ values to 4 decimal places.

\num Sketch the graph of $y = x \ln{x}$ on the interval $0 < x < 1$.
Where does this function have its minimum, and what is the minimum
value?

\num Let $x$ be a positive number; if $x > 1$ then the cube of $x$ is
larger than its square.  However, if $0 < x < 1$ then the square is
larger than the cube.  How {\em much\/} larger can it be; that is,
what is the greatest amount by which $x^2$ can exceed $x^3$? For which
$x$ does this happen?

\num By how much can $x^p$ exceed $x^q$, when $0 < p < q$ and $0 < x$?
For which $x$ does this happen?

%%%%%%%%%%%%%%%%%%%%%%%%%%%%%%%%%%%%%%%%%%%%%%%%%%%%%%%%%%%%%%%%%%%%%%%%
%                                                                      %
%                               Section 4                              % 
%                                                                      %
%%%%%%%%%%%%%%%%%%%%%%%%%%%%%%%%%%%%%%%%%%%%%%%%%%%%%%%%%%%%%%%%%%%%%%%%


\section{Optimal Shapes}

\subsection{The Problem of the Optimal Tin Can}

%
\mar{\vspace{3pt}What is the minimum surface area?}      
%
\noindent
     \begin{minipage}{3.0in}
Suppose you are a tin can manufacturer.  You must make a can to hold a
certain volume $V$ of canned tomatoes.  Naturally, the can will be a
cylinder, but the proportions, the height $h$ and the radius $r$, can
vary. Your task is to choose the proportions so that you use the least
amount of tin to make the can.  In other words, you want the surface
area of that can to be as small as possible.
\end{minipage} \
     \begin{minipage}{1.75in}
     \vspace*{1.75in}
     \special{psfile=../ps/calc1/pic257a.ps}
     \end{minipage}

\subsection{The Solution}

     The surface area is the sum of the areas of the two circles at the top and bottom of the can, plus the area of the rectangle that would be obtained if the top and bottom were removed and the side cut vertically.
\vspace*{1.5in}

\special{psfile=../ps/calc1/pic257b.ps}

Thus $A$ depends on $r$ and $h$:
$$
A = 2 \pi r^2 + 2 \pi r h.
$$
However, $r$ and $h$ cannot vary independently.  
	\mar{Finding $A$ as a function of $r$}
Because the volume $V$ is fixed,  $r$ and $h$ are related by
$$
V = \pi r^2 h.
$$
Solving the equation above for $h$ in terms of $r$, 
$$
h = \frac{V}{\pi r^2},
$$
we may express $A$ as a function of $r$ alone:
$$
A = f(r) = 2 \pi r^2 + 2 \pi r \frac{V}{\pi r^2} = 2\pi r^2 + \frac{2V}{r}
$$
	\mar{$V$ is a parameter}
\noindent
Notice that the formula for this function involves the parameter $V\/$.  
The mathematical description of our task is to find the value of $r$ that makes $A=f(r)$ a minimum.

     Following the procedure of the previous section, we first determine the domain of the function.  
	\mar{Finding the domain\\ of $f(r)$}     
Clearly this problem makes physical sense only for $r>0$. Looking at the equation
$$
h = \frac{V}{\pi r^2},
$$
we see that although $V$ is fixed, $r$ can be arbitrarily large provided $h$ is sufficiently small (resulting in a can that looks like an elephant stepped on it).  Thus the domain of our function is $r>0$, which is not a closed interval, so we have no guarantee that a minimum exists.

   Next we compute $f'(r)$, keeping in mind that the symbols $V$ and $\pi$ represent constants and that we are differentiating with respect to the variable $r$.  
$$
f'(r) = 4 \pi r  - \frac{2V}{r^2} = \frac{4 \pi r^3 - 2V}{r^2}
$$
The derivative is undefined at $r=0$, which is outside the domain under consideration.  
	\mar{Looking for\\ critical points}
So now we set the derivative equal to zero and solve for any possible critical points. 
\begin{align*}
f'(r) &= \frac{4 \pi r^3 - 2V}{r^2}  \\[3pt]
    0 &= \frac{4 \pi r^3 - 2V}{r^2}  \\[3pt]
    0 &= 4 \pi r^3 - 2V  \\[3pt]
    r &= \sqrt[3]{V/2 \pi}
\end{align*}  
Thus $r=\sqrt[3]{V/2 \pi}$ is the only critical point.

We can actually sketch the shape of the graph of $A$ versus $r$ 
%
\mar{Finding the shape\\ of the graph of $A$}	
%
based on this analysis of $f'(r)$.  The sign of $f'(r)$ is determined
by its numerator, since the denominator $r^2$ is always positive.

\begin{itemize}
\item When $r < \sqrt[3]{V/2 \pi}$, the numerator is negative, so the graph of $A$ is falling.  
\item When 
$r>\sqrt[3]{V/2 \pi}$, the numerator is positive, so  the graph is rising.

\end{itemize}

\vspace{130pt}

\special{psfile=../ps/calc1/pic259.ps}

\noindent Obviously, $A$ has a {\em minimum\/} at 
$$
r_{\rm extreme}=\sqrt[3]{V/2 \pi}.
$$  
It is also obvious 
%
\mar{A has a minimum\\ but no maximum}
%
that there is no {\em maximum\/} area, since  
$$ 
\lim_{r \rightarrow 0} A = \infty \quad\mbox{and}\quad  \lim_{r \rightarrow \infty} A = \infty.
$$
Thus we see again that not {\em every\/} optimization problem has a solution.

\subsection{The Mathematical Context:  Optimal Shapes}

It is interesting to find the value of the height $h=h_{\rm extreme}$
when the area is a minimum.  Since
\[
h_\text{extreme} = \frac{V}{\pi r_\text{extreme}^2},
\]
we can use $r_\text{extreme}=\sqrt[3]{V/2 \pi}$ to write $V = 2 \pi
r_\text{extreme}^3$ and thus find
\[
h_\text{extreme} 
   = \dfrac{2 \pi r_\text{extreme}^3}{\pi r_\text{extreme}^2}
   = 2 r_\text{extreme}.
\]
In other words, the height of the optimal tin can exactly equals its
diameter.  Campbell soup cans are far from optimal, but a can of
Progresso plum tomatoes has diameter 4 inches and height 4.5 inches.
Does someone at Progresso know calculus?

Problems like the one
%                                                                              \mar{Graphing utilities are often best for solving practical problems}
%
above---as well as others that you will find in the exercises---have
been among the mainstays of calculus courses for generations.
However, if you really had to face the problem of optimizing the shape
of a tin can as a practical matter, you would probably have a
numerical value for $V$ specified in advance.  In that case, your
first impulse might be to use a graphing utility to approximate
$r_{\it extreme}$, as accurately as your needs warrant.  This is
entirely appropriate.  Using a graphing utility gives the answer as
quickly as taking derivatives, and you are less likely to make
mistakes in algebra.  Moreover, the functions encountered in many
other real problems are often complicated and messy, and their
derivatives are hard to analyze.  A graphing utility may provide the
only practical course open to you.


So why insist that you use derivatives on these problems?
%                                                                  
\mar{So why do things\\ the hard way?}
%
It is because the actual context of these examples is not really
saving money for food canners.  The real context is geometry.  If we
had used a graphing facility to solve the can problem, we might have
noticed that the optimal radius was about half the height, but we
wouldn't have known the relationship is exact.  Nor would we have
recognized that the relationship holds for cylinders of arbitrary
volume.

Using a graphing utility is essentially mindless, which is part of its
virtue if we just need a specific answer to a particular problem.
Part of the purpose of a calculus course, though, is to develop the
concepts, tools, and the geometric vocabulary for thinking about
problems in general, to see the connecting threads between apparently
disparate settings.  The more clearly we can think within a general
framework, not just that of the specific problem being addressed, the
better intuitions we will develop that will help us see unsuspected
relationships in the problem at hand or think more creatively about
other problems in other settings. The ability to give precise
expression to our intuitions can lead to deeper insights.  You should
thus try to see the exercises that follow as not just
%
\mar{Symmetric shapes\\ are often optimal}     
%
about fences and storage bins---they are opportunities to observe
that, often, the geometric regularity that pleases the eye is also
optimal.


\newpage

\subsection{Exercises}

For each of the following problems, find a function relating the
variables and use differentiation to find the optimal value specified.
Some of the problems (1, 3, 4, and 5) are expressed in terms of a
general parameter---$P$, $L$, $A$, and $L$ again---rather than
specific numerical values.  If this gives you trouble, try doing the
problem using the specific parameter value given at the end of the
problem for computer verification.  Use other specific values as
needed until you can do the problem in terms of the general parameter,
which should behave exactly like any of the specific values you tried.

\num Show that the rectangle of perimeter $P$ whose area is a maximum
is a square.  Use a graphing utility to check your answer for the
special case when $P$=100 feet.

\num An open rectangular box is to be made from a piece of cardboard 8
inches wide and 15 inches long by cutting a square from each corner
and bending up the sides.  Find the dimensions of the box of largest
volume.  Use a graphing utility to check your answer.

\num One side of an open field is bounded by a straight river.  A
farmer has $L$ feet of fencing.  How should the farmer proportion a
rectangular plot along the river in order to enclose as great an area
as possible?  Use a graphing utility to check your answer for the
special case when $L$ = 100 feet.

\num An open storage bin with a square base and vertical sides is to
be constructed from $A$ square feet of wood.  Determine the dimensions
of the bin if its volume is to be a maximum.  (Neglect the thickness
of the wood and any waste in construction.)  Use a graphing utility to
check your answer for the special case when $A$=100 square feet.

\num A roman window is shaped like a rectangle surmounted by a
semicircle.  If the perimeter of the window is $L$ feet, what are the
dimensions of the window of maximum area?  Use a graphing utility to
check your answer for the special case when $L$ = 100 feet.

\num Suppose the roman window of problem 5 has clear glass in its
rectangular part and colored glass in its semicircular part.  If the
colored glass transmits only half as much light per square foot as the
clear glass does, what are the dimensions of the window that transmits
the most light?  Use a graphing utility to check your answer for the
special case when $L$ = 100 feet.

\num A cylindrical oil can with radius $r$ inches and height $h$
inches is made with a steel top and bottom and cardboard sides.  The
steel costs 3 cents per square inch, the cardboard costs 1 cent per
square inch, and rolling the crimp around the top and bottom edges
costs 1/2 cent per linear inch.  (Both crimps are done at the same
time, so only count the contribution of one circumference.)  

\sub Express the cost $C$ of the can as a function of $r$ and $h\/$.

\sub Find the dimensions of the cheapest can holding 100 cubic inches
of oil.  (You'll need to solve a cubic equation to find the critical
point.  Use a graphing utility or an estimation procedure to
approximate the critical point to 3 decimal places.)

%%%%%%%%%%%%%%%%%%%%%%%%%%%%%%%%%%%%%%%%%%%%%%%%%%%%%%%%%%%%%%%%%%%%%%%%
%                                                                      %
%                               Section 5                              % 
%                                                                      %
%%%%%%%%%%%%%%%%%%%%%%%%%%%%%%%%%%%%%%%%%%%%%%%%%%%%%%%%%%%%%%%%%%%%%%%%

\section{Newton's Method}

\subsection{Finding Critical Points}

When we solve optimization problems for functions given by formulas,
we begin by calculating the derivative and using the derivative
formula to find critical points.  Almost always the derivative is
defined for all elements of the domain, and we find the critical
points by determining the roots of the equation obtained by setting
the derivative equal to zero.

Finding the roots of this equation often requires an estimation
procedure.  For example, consider the function
$$
f(x) = x^4 + x^3 + x^2 + x + 1.
$$
The derivative of $f$ is
$$
f'(x) = 4x^3 +3x^2 +2x + 1,
$$
which is certainly defined for all $x$.  

In order to use a graphing utility to find the roots of $f'(x)=0$, we
need to choose an interval that will contain the roots we seek.
%
\mar{Solving $f'(x)=0$}
%
Since $f'(0)=1>0$ and $f'(-1) = -2 < 0$, we know $f'$ has at least one
root on $[-1,0]$.  (Why is this?)  But might there be other roots
outside this interval?

It is easy to see that $f'(x)$ is positive for all $x>0$, so there are
no roots to the right of $[-1,0]$.  What about $x<-1$?  Rewriting the
derivative as
$$
f'(x) = (2x+1)(2x^2+1) + x^2
$$
lets us see that $f'(x)$ is negative for all $x<-1$ (check this for
yourself), so there are no roots to the left of $[-1,0]$ either.

Examining the graph of $y=4x^3 +3x^2 +2x + 1$ on $[-1,0]$, we see that
it crosses the $x$-axis exactly once, so there is a unique critical
point. Progressively shrinking the interval, we find that, to eight
decimal places, this critical point is
$$
c= -.60582958 \ldots .
$$
Notice that it in this method it requires roughly the same amount of
work to get each additional digit in the answer.

There is, however, another approximation procedure we can use, called
Newton's method.  This has wide applicability and usually converges
very rapidly to solutions---additional digits are much easier
%
\mar{A good algorithm is\\ a convenient tool}
%
to get than in the method we just looked at.  It also has the virtue
of being algorithmic, so that we can write a single computer program
which can then be used for any root-finding problem we may encounter
without a great deal of thought and decision-making.  Newton's method
is also an interesting application of both local linearity and
successive approximation, two important themes of this course.

\subsection{Local Linearity and the Tangent Line}
 

Let's give the derivative function the new name $g$ to emphasize that
we now want to consider it as a function in its own right,
%
\mar{The root is an $x$-intercept of the graph}     
%
out of the context of the function $f$ from which it was derived.
Look at the following graph of the function
	$$
	y = g(x) = 4x^3 +3x^2 + 2x + 1\,.
	$$
We are seeking the number $r$ such that $g(r) = 0$.  The root $r$ is
the $x$-coordinate of the point where the graph crosses the $x$-axis.

The basic plan of attack in Newton's method is to replace the graph of
$y=g(x)$ by a straight line that looks reasonably like the graph near
the root $r$.Then, the $x$-intercept of that line will be a reasonably
good estimate for $r$.  The graph below includes such a line.

\mbox{}

\vspace{140pt}

\special{psfile=../ps/calc1/pic280a.ps}

\special{psfile=../ps/calc1/pic280b.ps}

The function $g$ is locally linear at $x=-1\/$, and we have drawn an
extension of the local linear approximation of $g$ at this point.
%
\mar{The tangent line extends the local linear approximation}
%
This line is called the {\bf tangent line\/}\index{tangent line} to
the graph of $y=g(x)$ at $x=-1\/$, by analogy to the tangent line to a
circle.  To find the $x$-intercept of the tangent line, we must know
its equation.  Clearly the line passes through the {\bf point of
  tangency}\index{tangency!point of} $(-1,g(-1))=(-1,-2)$.  What is
its slope?  It is the same as the slope of the local linear
approximation at $(-1,g(-1))$, namely
$$
g'(-1)=12(-1)^2 + 6(-1) + 2 = 8.
$$
Thus the equation
%
\mar{The equation of\\ the tangent line}
%
of the tangent line is
$$
y + 2 = 8(x+1).
$$

To find the $x$-intercept of this line, we must set $y$ equal to zero and solve for $x$: \  $0+2=8(x+1)$  gives us  $x=-0.75$.  
%
\mar{Finding the $x$-intercept of the tangent line}     
%
Of course, this $x$-intercept is not equal to $r$, but it's a better
approximation than, say, --1.  To get an even better approximation, we
repeat this process, starting with the line tangent to the graph of
$g$ at $x=-0.75$ instead of at $x=-1$.

\vspace*{150pt}

\special{psfile=../ps/calc1/pic281.ps}

The slope of this new tangent line is $g'(-0.75) = 4.25$, and it
passes through $(-0.75, -0.5)$, so its equation is
$$
y + 0.5 = 4.25(x+0.75).
$$ 
Setting $y=0$ and solving for $x$ gives a new $x$-intercept whose
value is equal to $-0.6323529$, closer still to $r$.

It seems reasonable to repeat the process yet again, using the tangent
line at $x = -0.6323529$. Let's first introduce some notation to keep
track of our computations.
%
\mar{Successive approximations get closer to the root $r$}    
%
Call the original point of tangency $x_0$, so $x_0=-1$.  Let $x_1$ be
the $x$-intercept of the tangent line at $x=x_0$.  Draw the tangent
line at $x=x_1$, and call its $x$-intercept $x_2$.  Continuing in this
way, we get a sequence of points $x_{0}, x_{1}, x_{2}, x_{3}, \ldots$
which appear to approach nearer and nearer to $r$.  That is, they
appear to approach $r$ as their limit.

\subsection{The Algorithm}

This process of using one number to determine the next number in the
sequence is the heart of Newton's method---it is an iterative method.
Moreover, it turns out to be quite simple to calculate each new
estimate in terms of the previous one.  To see how this works, let's
compute $x_{1}$ in terms of $x_{0}$.  We know $x_{1}$ is the
$x$-intercept of the line tangent to the graph of $g$ at $(x_{0},
g(x_{0}))$.  The slope of this line is $g'(x_{0})$,
%
\mar{\vspace{12pt}The general equation of the tangent line} 
%
so 
	$$  
y - g(x_{0}) = g'(x_{0})(x - x_{0}) 
	$$ 
is the {\bf equation of the tangent line}\index{tangent line}.  Since this line crosses the $x$-axis at the point  $(x_{1},0)$,  we set  $x = x_{1}$  and  $y = 0$  in the equation to obtain
	$$
 0 - g(x_{0}) = g'(x_{0})(x_{1} - x_{0}).
	$$
Now it is easy to solve for  $x_{1}$:
\begin{align*}
     g'(x_{0})(x_{1} - x_{0}) &= -g(x_{0}) \\
                x_{1} - x_{0} &= \frac{-g(x_{0})}{g'(x_{0})}  \\
                        x_{1} &= x_{0} - \frac{g(x_{0})}{g'(x_{0})}.
\end{align*}
In the same way we get 
\[
x_{2} = x_{1} - \frac{g(x_{1})}{g'(x_{1})},  \qquad
x_{3} = x_{2} - \frac{g(x_{2})}{g'(x_{2})},       
\]
and so on.

To summarize, suppose that $x_{0}$ is given some value START.  Then
Newton's method\index{Newton's method} is the computation of the
sequence of numbers determined by
\begin{align*}
             x_{0} &= \mbox{START}    \\
           x_{n+1} &= x_{n} - \frac{g(x_{n})}{g'(x_{n})}, 
               \qquad n = 0,1,2,3,\ldots
\end{align*}
 
	As we have seen many times, the sequence
$$
x_1, x_2, x_3, \ldots, x_n, \ldots 
$$
is a list of numbers to which we can always add a new term---by iterating our method yet again.  
	\mar{The limit of\\ the successive approximations\\ is the root}
For most functions, if we begin with an appropriate starting value of $x_0\/$, there is another number $r$ that is the {\em limit\/} of this list of numbers, in the sense that the difference between $x_n$ and $r$ becomes as small as we wish as $n$ increases without bound,
	$$
	r= \lim_{n \rightarrow \infty} x_n.
	$$
The numbers $x_1, x_2, x_3, \ldots, x_n, \ldots$ constitute a sequence of
{\em successive approximations\/} for the root $r$ of the equation $g(x)=0\/$.
We can write a computer program to carry out this algorithm for as many steps as we choose.  The program NEWTON does just that for 
$g(x) = 4x^3 + 3x^2 +2x +1\/$.

\begin{center}
{\bf Program:  NEWTON}
\index{program!NEWTON}

{\bf Newton's method for solving $g(x)=4x^3 + 3x^2 +2x +1 = 0$}
\end{center}
\begin{verbatim}
   start = -1
   numberofsteps = 8
   x = start    
   FOR n = 0 to numberofsteps   
      print n, x
      g = 4 * x^3 + 3 * x^2 + 2 * x + 1
      gprime = 12 * x^2 + 6 * x + 2
       x = x - g/gprime
    NEXT n   
\end{verbatim}


If we program a computer using this algorithm with START = --1, then
we get
\begin{align*}
     x_0 &= -1.00000000000  \\
     x_1 &= -0.75000000000  \\
     x_2 &= -0.63235294118  \\
     x_3 &= -0.60687911790  \\
     x_4 &= -0.60583128240  \\
     x_5 &= -0.60582958619  \\
     x_6 &= -0.60582958619  \\
     x_7 &= -0.60582958619  \\
     x_8 &= -0.60582958619.
\end{align*}

Thus we have found the root of $g(x)=0$---the critical point we were
looking for.  In fact, after only 6 steps we could see that the value
of the critical point was specified to at least ten decimal places.
Also at the sixth step, we had the eight decimal places obtained with
the use of the graphing utility.  In fact, it turns out that the
number of decimal places fixed roughly doubles with each round.  In
the above list, for instance, $x_2$ fixed one decimal, $x_3$ fixed two
decimals, $x_4$ fixed four, $x_5$ fixed ten (the eleventh digit of the
root is really an 8, which gets rounded to a 9 in $x_6$ -- $x_8$.),
and $x_6$ would have fixed at least twenty if we had printed them all
out!  Moreover, by changing only three lines of the program---the
first, sixth, and seventh---we can use NEWTON for any other function.
With the use of the program NEWTON, we will see that in most cases we
can obtain results more quickly and to a higher degree of accuracy
with Newton's method than by using a graphing utility, although it is
still sometimes helpful to use a graphing utility to get a reasonable
starting value.


\subsection{Examples}

\noindent 
{\bf Example 1}.  Start with $\cos{x} = x$.  The solution(s) to this
equation (if any) will be the $x$-coordinates of any points of
intersection of the graphs of $y = \cos{x}$ and $y = x$.  Sketch
these two graphs and convince yourself that there is one solution,
between 0 and $\pi/2$.  The equation $\cos{x} = x$ is not in the
form $g(x) = 0$, so rewrite it as $\cos{x} - x = 0$.  Now we can
apply Newton's method with $g(x) = \cos{x} - x$.  Try starting with
$x_{0} = 1$.  This gives the iteration scheme
\begin{align*}             
       x_{0} &= 1      \\
     x_{n+1} &= x_{n} - \frac{\cos{x_{n}} - x_{n}}{-\sin{x_{n}} - 1}, 
         \qquad n = 0,1,2, \ldots  
\end{align*}

 The numbers we get are
\begin{align*}
        x_{0} &= 1.000000000  \\
        x_{1} &= \hphantom{1}.750363868 \ldots \\
        x_{2} &= \hphantom{1}.739112891 \ldots \\
        x_{3} &= \hphantom{1}.739085133 \ldots \\
        x_{4} &= \hphantom{1}.739085133 \ldots 
\end{align*}
We have the solution to 9 decimal places in only four steps.  Not only
does Newton's method work, it works fast!

\bigskip
\noindent 
{\bf Example 2}.  Suppose we continue with the equation $\cos{x} = x$,
but this time choose $x_{0} = 0$.  What will we find?  The numbers we
get are
\begin{align*}
        x_{0} &= 0.000000000  \\
        x_{1} &= 1.000000000  \\
        x_{2} &= 0.750363868
\end{align*}   
There is no need to continue; we can see that we will again obtain $r
= 0.739085133 \ldots$ as in Example~1.  Look again at your sketch and
see why you might have predicted this result.

\bigskip
\noindent {\bf Example 3}.  Next, let's find the roots of the
polynomial $x^{5} - 3x + 1\/$.  This means solving the equation $x^{5}
- 3x + 1 = 0\/$.  We know the necessary derivative,
%
\mar{Finding the starting value $x_0$ can be hard} 
%
so we're ready to apply Newton's method, except for one thing: which
starting value $x_{0}$ do we pick?  This is the part of Newton's
method that leaves us on our own.

Assuming that some graphing software is available, the best thing to
do is graph the function.  But for most graphing utilities, we need to
choose an interval.  How do we choose one which is sure to include all
the roots of the polynomial?  The derivative $5x^4 -3$ of this
polynomial is simple enough
%
\pagebreak
%
that we can use it to get an idea of the shape of the graph of
$y=x^5-3x+1$ before we turn to the computer.  Clearly the derivative
is zero only for $x = \pm \sqrt[4]{3/5}\/$, and the derivative is
positive except between these two values of $x\/$.
%
\mar{Using the derivative\\ to find the shape\\ of the graph} 
%
In other words, we know the shape of the graph of $y=g(x)\/$---it is
increasing, then decreases for a bit, then increases from there on
out.  This still is not enough information to tell us how many roots
$g$ has, though; its graph might lie in any one of the following
configurations and so have 1, 2, or 3 roots (there are two other
possibilities not shown---one has 2 roots and one has 1 root).
%
\mar{\special{psfile=../ps/calc1/pic285.ps}}
%
We can thus say that $g(x)=0$ has at least one and at most three real
roots.  However, if we further observe that $g(-2) = -25, g(-1)=3,
g(0) = 1,g(1)= -1$, and $g(2)=27$, we see that the graph of $g$ must
cross the $x$--axis at some value of $x$ between -2 and -1, between 0
and 1, and again between 1 and 2.  Therefore the right--hand sketch
above must be the correct one.

Or we can almost as easily turn to a graphing utility.  If we try the
interval $[-5,5]\/$, we see again that the graph crosses the $x$-axis
in exactly three points.

\vspace{135pt}
  
\special{psfile=../ps/calc1/pic286.ps}

\enlargethispage*{20pt}
One of the roots is between $-2$ and $-1$, 
%
\mar{Finding the root between \\$-2$ and $-1$}
%
one is between 0 and 1, and the third is between 1 and 2.  To find the
first, we apply Newton's method with $x_{0} = -2\/$.  Then we get
\begin{align*}
              x_{0} &= -2.000000000  \\
              x_{1} &= -1.675324675 \ldots  \\
              x_{2} &= -1.478238029 \ldots  \\
              x_{3} &= -1.400445373 \ldots  \\
              x_{4} &= -1.389019863 \ldots  \\
              x_{5} &= -1.388792073 \ldots  \\
              x_{6} &= -1.38879198 4\ldots  \\
              x_{7} &= -1.388791984 \ldots   
\end{align*}
This took a few more steps than the other examples, but not a lot.
Notice again that once there are any decimals fixed at all, the number
of decimals fixed roughly doubles in the next approximation. In the
exercises you will be asked to compute the other two roots.

\bigskip
\noindent {\bf Example 4}.  Let's use Newton's method to find the
obvious solution $r = 0$ of $x^{3} - 5x = 0\/$.  If we choose $x_{0}$
sufficiently close to 0, Newton's method should work just fine.  But
what does ``sufficiently close'' mean?  Suppose we try $x_{0} = 1\/$.
Then we get
\begin{align*}       
            x_{0} &= \hphantom{-}1  \\
            x_{1} &= -1 \\
            x_{2} &= \hphantom{-}1  \\
            x_{3} &= -1 \\
            x_{4} &= \hphantom{-}1  \\
                  &\ \vdots   
\end{align*}
The  $x_{n}\/$'s oscillate endlessly, never getting close to  0.  
%
\mar{Newton's method\\ can fail}
%
Going back to the geometric interpretation of Newton's method, this
oscillation can be explained by the graph below.

\vspace{145pt}

\special{psfile=../ps/calc1/pic287.ps}

Using more advanced methods, it is possible to get precise estimates
for how close $x_{0}$ needs to be to $r$ in order for Newton's method
to succeed.  For now, we'll just have to rely on common sense and
trial and error.

One important thing to note is the relation between algebra and
Newton's method.  Although we can now solve many more equations than
we could earlier, this doesn't mean that we can abandon algebra.  In
fact, given a new equation, you should first try to solve it
algebraically, for exact solutions are often better.  Only when this
fails should you look for approximate solutions using Newton's method.
So don't forget algebra---you'll still need it!


\subsection{Exercises}

\num {\bf The Babylonian algorithm.}\index{Babylonian algorithm} Show
that the Babylonian algorithm of chapter 2 is the same as Newton's
method applied to the equation $x^2-a=0\/$.

\num When Newton introduced his method, he did so with the example
$x^3 - 2x - 5 = 0$.  Show that this equation has only one root, and
find it.

\aside{This example appeared in 1669 in an unpublished manuscript of
  Newton's (a published version came later, in 1711).  The interesting
  fact is that Newton's method {\em differs\/} from the one presented
  here: his scheme was more complicated, requiring a different formula
  to get each approximation.  In 1690, Joseph Raphson\index{Raphson,
    Joseph} transformed Newton's scheme into the one used above.
  Thus, ``Newton's method'' is more properly called the
  ``Newton--Raphson method'', and many modern texts use this more
  accurate name.}

\num Use Newton's method to find a solution of $x^3 + 2x^2 + 10x = 20$
near the point $x=1$.

\aside{The approximate solution 1;22,7,42,33,4,40 of this equation
  appears in a book written in 1228 by Leonardo of Pisa\index{Leonardo
    of Pisa} (also known as \mbox{Fibonacci}\index{Fibonacci}).  This
  number looks odd because it's written in sexagesimal notation: it
  translates into
$$
1+\frac{22}{60} +\frac{7}{60^2}+\frac{42}{60^3}+\frac{33}{60^4}+\frac{4}{60^5}+\frac{40}{60^6}.
$$
This solution is accurate to 10 decimal places, which is not bad for
750 years ago.  In the Middle Ages, there was a lot of interest in
solving equations. There were even contests, with a prize going to the
person who could solve the most. The quadratic formula, which
expresses algebraically the roots of any second degree equation, had
been known for thousands of years, but there were no general methods
for finding roots of higher degree equations in the 13th century.  We
don't know how Leonardo found his solution---why give away your
secrets to your competitors!}

\num  Use Newton's method to find a solution of $x^3+3x^2=5$.

\aside{In 1530, Nicolo Tartaglia\index{Tartaglia, Nicolo} was
  challenged to solve this equation algebraically.  Five years later,
  in 1535, he found the solution
$$
x = \sqrt[3]{\frac{3+\sqrt{5}}{2}}+\sqrt[3]{\frac{3-\sqrt{5}}{2}}-1 \, .
$$
Initially, Tartaglia could solve only certain types of cubic
equations, but this was enough to let him win some famous contests
with other mathematicians of the time.  By 1541, he knew the general
solution, but he made the mistake of telling Geronimo
Cardano\index{Cardano, Geronimo}.  Cardano published the solution in
1545 and the resulting formulas are called ``Cardan's
Formulas''.\index{Cardan's formulas}}

The above solution of $x^3+3x^2=5$ is called a {\bf solution by
  radicals}\index{solution!by radicals} because it is obtained by
extracting various roots or radicals.  Similarly, some time before
1545, Luigi Ferrari\index{Ferrari, Luigi} showed that any fourth
degree equation can be solved by radicals.  This led to an intense
interest in the fifth degree equation.  To see what happens in this
case, read the next problem.

\num In Example 3, we saw that one root of $x^5-3x+1$ was
--1.39887919\ldots . Use Newton's method to find the other two roots.

\aside{In 1826, Niels Henrik Abel\index{Abel, Niels} proved that the
  general polynomial of degree 5 or greater cannot be solved by
  radicals.  Using the work of Evariste Galois\index{Galois, Evariste}
  (done around 1830, but not understood until many years later), it
  can be shown that the roots of the equation $x^5-3x+1=0$ cannot be
  expressed by any combination of radicals.  Thus algebra can't solve
  this equation---some kind of successive approximation technique is
  unavoidable!}

\num One of the more surprising applications of Newton's method is to
compute reciprocals.  To make things more concrete, we will compute
1/3.4567 .  Note that this number is the root of the equation $1/x =
3.4567$.

\sub Show that the formula of Newton's method gives us
\[
x_{n+1} = 2 x_n-3.4567 x_n^2.
\]

\sub Using $x_0=.5$ and the formula from (a), compute 1/3.4567 to a
high degree of accuracy.

\sub Try starting with $x_0=1$.  What happens? Explain graphically
what goes wrong.

This method for computing reciprocals is important because it involves
only {\em multiplication\/} and {\em subtraction\/}.  Since
$a/b=a\cdot (1/b)$, this implies that division can likewise be built
from multiplication and subtraction.  Thus, when designing a computer,
the division routine doesn't need to be built from scratch---the
designer can use the method illustrated here.  There are some
computers that do division this way.

\num In this problem we will determine the maximum value of the function
\[
f(x) = \dfrac{x+1}{x^4+1}.
\]

\sub Graph $f(x)$ and convince yourself that the maximum value occurs
somewhere around $x=.5$.  Of course, the exact location is where the
slope of the graph is zero, i.e., where $f'(x)=0$.  So we need to
solve this equation.

\sub Compute $f'(x)$.

\sub Since the answer to (b) is a fraction, it vanishes when its
numerator does.  Setting the numerator equal to 0 gives a fourth
degree equation.  Use Newton's method to find a solution near $x=.5 \,
.$

\sub Compute the maximum value of $f(x)$.

\num  Consider the hyperbola $y=1/x$ and the circle $x^2-4x+y^2+3=0$.

\sub By graphing the circle and the hyperbola, convince yourself that
there are two points of intersection.

\sub By substituting $y=1/x$ into the equation of the circle, obtain a
fourth degree equation satisfied by the $x$--coordinate of the points
of intersection.

\sub Solve the equation from (b) by Newton's method, and then
determine the points of intersection.

\num Sometimes Newton's method doesn't work so nicely.  For example,
consider the equation $\sin x = 0$.

\sub Compute $x_1$ using Newton's method for each of the four starting
values $x_0=1.55,1.56,1.57 \mbox{ and } 1.58$.

\sub The answers you get are wildly different.  Using the basic
formula
\[
x_{n+1}=x_n-\frac{g(x_n)}{g'(x_n)}
\]
explain why.




\subsubsection{The epidemic runs its course}

We return to the epidemiology example we have studied since chapter~1.
Recall that our $S$-$I$-$R$ model keeps track of three subgroups of
the population: the susceptible, the infected, and the recovered.  One
of the interesting features of the model is that the larger the
initial susceptible population, the more rapidly the epidemic runs its
course.  We observe this by choosing fixed values of $R_0 = R(0)$ and
$I_0 = I(0)$ and looking at graphs of $S(t)$ versus $t$ for various
values of $S_0 = S(0)$.

\newpage

\mbox{}

\vspace{170pt}

\special{psfile=../ps/calc1/pic291.ps}

\noindent
We see in each case that for sufficiently large $t$ the graph of $S$
levels off, approaching a value we'll call $S_\infty$:
\[
S_\infty = \lim_{t \rightarrow \infty} S(t).
\]

What we mean by the epidemic ``running its course'' is that $S(t)$
reaches this limit value.  We can see from the graphs that the value
of $S_0$ affects the number $S_\infty$ of individuals who escape the
disease entirely. It turns out that we can actually find the value of
$S_\infty$ if we know the values of $S_0, I_0,$ and the parameters $a$
and $b$.

Recall that $a$ is the {\em transmission coefficient}, and $b$ is the
{\em recovery coefficient\/} for the disease. The differential
equations of the $S$-$I$-$R$ model are
\begin{align*}
	S' &= -aSI, \\
	I' &= aSI - bI, \\
	R' &= bI. 
\end{align*}

\num Use the differentiation rules together with these differential
equations to show that
\[
(I+S-(b/a)\cdot \ln S)'=0.
\]

\num Explain why the result of problem 10 means that $I+S-(b/a)\cdot
\ln S$ has the same value---call it $C$---for every value of $t$.

\newpage

\num Look at the graphs of the solutions $I(t)$ for various values of
$S(0)$ below.

\vspace{200pt}

\special{psfile=../ps/calc1/pic292.ps}

\noindent
Write $\lim_{t \rightarrow \infty} I(t) = I_\infty$ .  What is the
value of $I_\infty$ for all values of $S(0)$?

\num Use the results of problems 11 and 12 to explain why
\[
S_\infty-\frac{b}{a}\ln (S_\infty)=I_0+S_0-\frac{b}{a}\ln (S_0).
\]
This equation determines $S_\infty$ implicitly as a function of $I_0$
and $S_0$.  For particular values of $I_0$ and $S_0$ (and of the
parameters), you can use Newton's method to find $S_\infty$.

\num  Use the values
\begin{align*}
a & =  .00001 \mbox{ (person-days)}^{-1}\\
b & = .08 \mbox{ day}^{-1}\\
I_0 &=  100 \mbox{ persons}\\
S_0 & =  35,000 \mbox{ persons}
\end{align*}
Writing $x$ instead of $S_\infty$ gives
\[
x-8000\ln (x) = -48,605.
\]
Apply Newton's method to find $x = S_\infty$.  Judging from the graph
for $S_0=35000$, it looks like a reasonable first estimate for
$S_\infty$ might be 100.

\num  Using the same values of $a$, $b$, and $I_0$ as in problem 14, determine the value of $S_\infty$ for each of the following initial population sizes:
\sub $S_0=45,000$.
\sub $S_0=25,000$.
\sub $S_0=5,000$.

%%%%%%%%%%%%%%%%%%%%%%%%%%%%%%%%%%%%%%%%%%%%%%%%%%%%%%%%%%%%%%%%%%%%%%%%
%                                                                      %
%                               Section 6                              % 
%                                                                      %
%%%%%%%%%%%%%%%%%%%%%%%%%%%%%%%%%%%%%%%%%%%%%%%%%%%%%%%%%%%%%%%%%%%%%%%%

\section{Chapter Summary}

\subsection{The Main Ideas}

\begin{itemize}

\item {\bf Formulas} for {\bf derivatives} can be calculated for
  functions given by formulas using the definition of the derivative
  as the limit of difference quotients
$$
\lim_{\Delta x \rightarrow 0} \frac{f(x+\Delta x)-f(x)}{\Delta x}.
$$

\item Each particular difference quotient is an {\bf approximation} to
  the derivative.  Successive approximations, for smaller and smaller
  $\Delta x\/$, approach the derivative as a limit.  The formula for
  the derivative gives its {\em exact\/} value.

\item Formulas for {\bf partial derivatives} are obtained using the
  formulas for derivatives of functions of a single variable by simply
  treating all variables other than the one of interest as if they
  were constants.

\item {\bf Optimization} problems occur in many different contexts.
  For instance, we seek to maximize benefits, minimize energy, and
  minimize error.

\item The {\bf sign} of the derivative indicates where a graph rises
  and where it falls.

\item Functions which are {\bf continuous} on a {\bf finite closed
  interval} have {\bf global extremes}.

\item For a function continuous on an interval, its {\bf local
  extremes} occur at {\bf critical points}---points where the
  derivative equals zero or fails to exist---or at endpoints.

\item Local linearity permits us to replace the graph of a function
  $y=g(x)$ by its {\bf tangent line} at a point near a root of
  $g(x)=0\/$, and then the $x$-intercept of the tangent line is a
  better approximation to the root.  Successive approximations,
  obtained by iterating this procedure, yield {\bf Newton's method}
  for solving the equation $g(x)=0\/$.

\end{itemize}

\subsection{Expectations}

\begin{itemize}

\item You should be able to {\bf differentiate} a function given by a
  formula.

\item You should be able to use differentiation formulas to {\bf
  calculate} partial derivatives.

\item In most cases, you should be able to determine from a graph of a
  function on an interval whether that function is {\bf continuous}
  and/or {\bf differentiable} on that interval.

\item You should be able to find {\bf critical points} for a function
  of one or several variables given by a formula.

\item You should be able to use the formula for the derivative of a
  function of a single variable to find {\bf local and global
    extremes}.

\item You should be able to use {\bf Newton's method} to solve an
  equation of the form $g(x)=0$.

\end{itemize}

\subsection{Chapter Exercises}

\subsubsection{Prices, demand and profit}  

Suppose the demand\index{demand} $D$ (in units sold) for a particular
product is determined by its price $p$ (in dollars), $D=f(p)$.  It is
reasonable to assume that when the price is low, the demand will be
high, but as the price rises, the demand will fall.  In other words,
we assume that the slope of the demand function is negative.  If the
manufacturing cost for each unit of the product is $c$ dollars, then
the profit per unit at price $p$ is $p-c$.  Finally the total profit
$T$ gained at the unit price $p$ will be the number of units sold at
the price $p$ (that is, the demand $D(p)$) multiplied by the profit
per unit $p-c$.
	$$
	T = g(p) = D(p)\, \mbox{units} \times (p-c) \, 
		\frac{\mbox{dollars}}{\mbox{unit}}
	$$
In this series of problems we will determine the effect of the demand
function and of the unit manufacturing cost on the maximum total
profit.

\num Suppose the demand function is linear
$$
D = f(p) = 1000 - 500p \  \mbox{units},
$$
and the unit manufacturing cost is .20, so the total profit is
$$
T = g(p) = (1000 - 500p)(p-.20) \  \mbox{dollars}.
$$
Find the ``best" price -- that is, find the price that yields the
maximum total profit.

\num Suppose the demand function is the same as in problem 1, but the
unit manufacturing cost rises to .30.  What is the ``best" price now?
How much of the rise in the unit manufacturing cost is passed on to
the consumer if the manufacturer charges this best price?

\num Suppose the demand function for a particular product is $D(p)=
2000-500p$, and that the unit manufacturing cost is .30.  What price
should the manufacturer charge to maximize her profit?  Suppose the
unit manufacturing cost rises to .50.  What price should she charge to
maximize her profit now?  How much of the rise in the unit
manufacturing cost should she pass on to the consumer?

\num If the demand function for a product is $D(p)=1500-100p$, compare
the ``best" price for unit manufacturing costs of .30 and .50.  How
much of the rise in cost should the manufacturer pass on to the
consumer?

\num As you may have noticed, problems 2--4 illustrate an interesting
phenomenon.  In each case, exactly half of the rise in the unit
manufacturing cost should be passed on to the consumer.  Is this a
coincidence?

\sub  Consider the most general case of a linear demand function
$$
D=f(p) = a-mp.
$$
and unit cost $c$.  What is the ``best" price?  Is exactly half the
unit manufacturing cost passed on to the consumer?  Explain your
answer.

\sub  Now consider a non-linear demand function
$$
D = f(p) = \frac{1000}{1+p^2}.
$$
Find the ``best" price for unit costs

	i) .50 dollar per unit;

	ii) 1.00 dollar per unit;

	iii) 1.50 dollars per unit.

\noindent 
How much of the price increase is passed on to the consumer in cases
(ii) and~(iii)?


